\documentclass[11pt,a4paper]{article}

% ==================== TEMEL PAKETLER (xelatex) ====================
\usepackage{fontspec}
\usepackage{polyglossia}
\setmainlanguage{turkish}
\usepackage[a4paper,margin=2cm,top=2.4cm,bottom=2.4cm]{geometry}
\usepackage{xcolor}
\usepackage{graphicx}
\usepackage{array}
\usepackage{booktabs}
\usepackage{longtable}
\usepackage{tabularx}
\usepackage{colortbl}
\usepackage{multirow}
\usepackage{enumitem}
\usepackage[most]{tcolorbox}
\usepackage{listings}
\usepackage{fancyhdr}
\usepackage{titlesec}
\usepackage{titletoc}
\usepackage{hyperref}
\usepackage{fontawesome5}

% ==================== FONTLAR ====================
\setmainfont{Latin Modern Roman}
\newfontfamily\codefont{Consolas}[Scale=0.90]

% ==================== RENK PALETİ ====================
\definecolor{anaMavi}{HTML}{1B3A5B}
\definecolor{vurguMavi}{HTML}{2E6DA4}
\definecolor{acikMavi}{HTML}{EAF2F9}
\definecolor{turuncu}{HTML}{D9701E}
\definecolor{yesil}{HTML}{2E7D32}
\definecolor{acikYesil}{HTML}{E8F5E9}
\definecolor{kirmizi}{HTML}{C0392B}
\definecolor{mor}{HTML}{6C3483}
\definecolor{acikMor}{HTML}{F3E9F7}
\definecolor{acikGri}{HTML}{F4F5F7}
\definecolor{ortaGri}{HTML}{6B7280}
\definecolor{koyuGri}{HTML}{2D2D2D}
\definecolor{kodArka}{HTML}{FBFBFA}
\definecolor{satirRenk}{HTML}{F0F4F8}

% ==================== HYPERREF ====================
\hypersetup{
  colorlinks=true,
  linkcolor=vurguMavi,
  urlcolor=vurguMavi,
  pdftitle={x86 Assembly Register ve BIOS Interrupt Rehberi},
  pdfauthor={Assembly Rehberi},
  bookmarksnumbered=true
}

% ==================== BAŞLIK STİLLERİ ====================
\titleformat{\section}
  {\color{anaMavi}\Large\bfseries}
  {\thesection.}{0.6em}{}
  [\vspace{2pt}{\color{turuncu}\titlerule[1.6pt]}]
\titleformat{\subsection}
  {\color{vurguMavi}\large\bfseries}
  {\thesubsection}{0.6em}{}
\titleformat{\subsubsection}
  {\color{koyuGri}\normalsize\bfseries}
  {\thesubsubsection}{0.6em}{}
\titlespacing*{\section}{0pt}{16pt}{8pt}

% ==================== BAŞLIK / ALTBİLGİ ====================
\pagestyle{fancy}
\fancyhf{}
\renewcommand{\headrulewidth}{0.4pt}
\renewcommand{\footrulewidth}{0.4pt}
\fancyhead[L]{\small\color{ortaGri}x86 Assembly Rehberi}
\fancyhead[R]{\small\color{ortaGri}\leftmark}
\fancyfoot[C]{\small\color{ortaGri}\thepage}

% ==================== KOD LİSTELEME ====================
\lstdefinestyle{asmstyle}{
  language={[x86masm]Assembler},
  basicstyle=\codefont\small\color{koyuGri},
  keywordstyle=\color{vurguMavi}\bfseries,
  commentstyle=\color{yesil}\itshape,
  stringstyle=\color{turuncu},
  numberstyle=\tiny\color{ortaGri},
  numbers=left,
  numbersep=8pt,
  backgroundcolor=\color{kodArka},
  frame=single,
  framesep=6pt,
  rulecolor=\color{acikGri!80!ortaGri},
  breaklines=true,
  showstringspaces=false,
  tabsize=4,
  columns=fullflexible,
  keepspaces=true,
  morekeywords={mov,int,org,bits,hlt,jmp,cli,sti,lea,push,pop,call,ret,add,sub,inc,dec,cmp,je,jne,jz,jnz,jl,jg,jle,jge,ja,jb,jae,jbe,jc,jnc,loop,xor,and,or,not,neg,shl,shr,sal,sar,rol,ror,in,out,rep,repe,repne,stosb,stosw,lodsb,lodsw,movsb,movsw,scasb,cld,std,db,dw,dd,dq,times,equ,section,global,syscall,mul,imul,div,idiv,test,lgdt,lidt,movzx,movsx,pusha,popa,iret,nop,resb,resw,cdq,addps,movups}
}
\lstset{style=asmstyle}

% ==================== ÖZEL KUTULAR ====================
\newtcolorbox{bilgikutu}[1][]{
  colback=acikMavi, colframe=vurguMavi, boxrule=0.8pt, arc=3pt,
  left=8pt,right=8pt,top=6pt,bottom=6pt,
  fonttitle=\bfseries\color{white}, coltitle=white,
  title={\faInfoCircle~~Bilgi}, #1
}
\newtcolorbox{ipucukutu}[1][]{
  colback=acikYesil, colframe=yesil, boxrule=0.8pt, arc=3pt,
  left=8pt,right=8pt,top=6pt,bottom=6pt,
  fonttitle=\bfseries\color{white}, coltitle=white,
  title={\faLightbulb~~İpucu}, #1
}
\newtcolorbox{uyarikutu}[1][]{
  colback=kirmizi!8, colframe=kirmizi, boxrule=0.8pt, arc=3pt,
  left=8pt,right=8pt,top=6pt,bottom=6pt,
  fonttitle=\bfseries\color{white}, coltitle=white,
  title={\faExclamationTriangle~~Dikkat}, #1
}
\newtcolorbox{ornekkutu}[1][Örnek]{
  colback=acikGri, colframe=ortaGri!60, boxrule=0.7pt, arc=2pt,
  left=8pt,right=8pt,top=6pt,bottom=6pt,
  fonttitle=\bfseries\color{koyuGri}, coltitle=koyuGri,
  title={\faCode~~#1}
}
\newtcolorbox{notkutu}[1][Not]{
  colback=acikMor, colframe=mor, boxrule=0.8pt, arc=3pt,
  left=8pt,right=8pt,top=6pt,bottom=6pt,
  fonttitle=\bfseries\color{white}, coltitle=white,
  title={\faStickyNote~~#1}
}

% Yardımcı komutlar
\newcommand{\reg}[1]{\texttt{\color{anaMavi}\textbf{#1}}}
\newcommand{\hx}[1]{\texttt{\color{turuncu}#1}}
\newcommand{\kod}[1]{\texttt{\color{koyuGri}#1}}

\renewcommand{\arraystretch}{1.35}

% ==================== BAŞLANGIÇ ====================
\begin{document}

% ---------- KAPAK ----------
\begin{titlepage}
\centering
\vspace*{1.6cm}
{\color{turuncu}\rule{\textwidth}{2pt}}\\[6pt]
{\color{anaMavi}\Huge\bfseries x86 Assembly}\\[10pt]
{\color{vurguMavi}\LARGE Register \& BIOS Interrupt Rehberi}\\[8pt]
{\color{ortaGri}\large 16-bit \textbullet\ 32-bit \textbullet\ 64-bit Kapsamlı Cheat Sheet}\\[6pt]
{\color{turuncu}\rule{\textwidth}{2pt}}\\[1.6cm]

\begin{tcolorbox}[width=0.88\textwidth,colback=acikMavi,colframe=vurguMavi,arc=4pt,boxrule=1pt]
\centering\large
Bu belge x86 mimarisindeki tüm register (yazmaç) ailelerini, adresleme
modlarını, bayrakları, SIMD/FPU birimlerini ve en sık kullanılan BIOS
kesmelerini \textbf{çok sayıda çalışan kod örneğiyle} detaylı biçimde açıklar.
\end{tcolorbox}

\vspace{0.9cm}
\begin{tcolorbox}[width=0.88\textwidth,colback=acikGri,colframe=ortaGri!50,arc=3pt,boxrule=0.6pt]
\small
\textbf{Kapsam:} Genel amaçlı register'lar \textbullet\ adresleme modları \textbullet\
segment/işaretçi register'ları \textbullet\ FLAGS \& koşullu atlamalar \textbullet\
MMX/SSE/AVX/x87 \textbullet\ kontrol register'ları \textbullet\ komut referansı \textbullet\
INT 10h/13h/15h/16h/1Ah \textbullet\ tam iki aşamalı önyükleyici \textbullet\
korumalı moda geçiş
\end{tcolorbox}

\vfill
{\color{ortaGri}\large Hazırlanma Tarihi: 9 Temmuz 2026}\\[4pt]
{\color{ortaGri} Gerçek Mod / Korumalı Mod / Uzun Mod}
\vspace*{0.6cm}
\end{titlepage}

% ---------- İÇİNDEKİLER ----------
\tableofcontents
\thispagestyle{empty}
\clearpage

% =====================================================================
\section{Giriş: x86 Mimarisine Genel Bakış}
% =====================================================================

x86, Intel'in 1978'de çıkardığı 8086 işlemcisinden gelişerek bugüne ulaşmış
bir işlemci mimarisidir. Mimari, geriye dönük uyumluluğu (backward
compatibility) sayesinde 45 yıl önceki 16-bit kodları hâlâ çalıştırabilir.
Register'lar da bu evrimi yansıtır: her yeni nesil, öncekinin register'larını
\emph{genişleterek} üzerine ekleme yapmıştır.

\begin{bilgikutu}
\textbf{Register (Yazmaç) nedir?} İşlemcinin \emph{içinde} bulunan, RAM'den çok
daha hızlı erişilen küçük hafıza hücreleridir. İşlemci aritmetik, mantık ve veri
taşıma işlemlerini doğrudan register'lar üzerinde yapar. RAM erişimi onlarca
nanosaniye sürerken register erişimi tek bir saat çevriminin küçük bir
kısmında gerçekleşir. Bu yüzden performans kritik kodlarda "veriyi
register'da tutmak" temel bir prensiptir.
\end{bilgikutu}

\subsection{İşlemci Çalışma Modları}

\begin{center}
\begin{tabularx}{\textwidth}{@{}>{\bfseries\color{anaMavi}}l >{\raggedright\arraybackslash}X c@{}}
\toprule
\rowcolor{acikMavi} \color{anaMavi}\textbf{Mod} & \color{anaMavi}\textbf{Açıklama} & \color{anaMavi}\textbf{Genişlik} \\
\midrule
Gerçek Mod & 8086 uyumlu, segment:offset adresleme, doğrudan donanım \& BIOS erişimi. Açılışta işlemci burada başlar. & 16-bit \\
\rowcolor{acikGri} Korumalı Mod & Sanal bellek, bellek koruması, koruma halkaları (ring 0--3), GDT/IDT tabloları & 32-bit \\
Uzun Mod & 64-bit, düz bellek modeli, geniş adres alanı, modern işletim sistemleri & 64-bit \\
\rowcolor{acikGri} Sanal 8086 & Korumalı mod içinde 16-bit programları çalıştırma alt modu & 16-bit \\
\bottomrule
\end{tabularx}
\end{center}

\begin{ipucukutu}
Bilgisayar açıldığında işlemci \textbf{daima gerçek modda (16-bit)} başlar.
BIOS kesmeleri yalnızca gerçek modda kullanılabilir. Önyükleyici (bootloader)
önce BIOS kesmeleriyle çekirdeği belleğe yükler, sonra korumalı/uzun moda
geçiş yapar. Geçişten sonra BIOS kesmeleri artık kullanılamaz.
\end{ipucukutu}

\subsection{Sözdizimi: Intel vs AT\&T}

Bu rehberde \textbf{Intel sözdizimi} (NASM tarzı) kullanılır çünkü daha okunaklıdır
ve önyükleyici geliştirmede standarttır. Farkı bilmek yararlıdır:

\begin{center}
\begin{tabularx}{\textwidth}{@{}>{\raggedright\arraybackslash}X >{\raggedright\arraybackslash}X@{}}
\toprule
\rowcolor{acikMavi}\color{anaMavi}\textbf{Intel (NASM)} & \color{anaMavi}\textbf{AT\&T (GAS)} \\
\midrule
\kod{mov eax, 5} & \kod{movl \$5, \%eax} \\
\rowcolor{acikGri} \kod{mov eax, [ebx]} & \kod{movl (\%ebx), \%eax} \\
Hedef solda: \kod{mov hedef, kaynak} & Hedef sağda: \kod{mov kaynak, hedef} \\
\bottomrule
\end{tabularx}
\end{center}

\begin{notkutu}[Intel sözdizimi kuralı]
\kod{komut hedef, kaynak} — yani \textbf{sol taraf hedeftir}. \kod{mov ax, bx}
"bx'i ax'e kopyala" demektir. Köşeli parantez \kod{[ ]} bellek erişimi anlamına
gelir: \kod{[ebx]} "ebx'in gösterdiği adresteki değer" demektir.
\end{notkutu}

% =====================================================================
\section{Genel Amaçlı Register'lar}
% =====================================================================

x86'da 64-bit modda 16 adet genel amaçlı register vardır. Bunların 8 tanesi
tarihsel köke sahiptir ve isimlerinin baş harfleri işlevlerinden gelir.

\subsection{Register Boyut Hiyerarşisi}

Aynı fiziksel register, farklı isimlerle farklı bit genişliklerinde
kullanılabilir. Örneğin \reg{RAX} 64-bit'in tamamıdır; alt 32 biti \reg{EAX},
alt 16 biti \reg{AX}, \reg{AX}'in üst baytı \reg{AH}, alt baytı \reg{AL}'dir.

\begin{center}
\begin{tabular}{|c|c|c|c|}
\hline
\rowcolor{anaMavi}\multicolumn{4}{|c|}{\color{white}\textbf{RAX -- 64 bit (bit 63 \dots 0)}}\\
\hline
\multicolumn{2}{|c|}{\cellcolor{acikGri}\textit{(üst 32 bit -- ayrı ismi yok)}} & \multicolumn{2}{c|}{\cellcolor{acikMavi}\textbf{EAX -- 32 bit}}\\
\hline
\multicolumn{3}{|c|}{\cellcolor{acikGri}} & \multicolumn{1}{c|}{\cellcolor{acikYesil}\textbf{AX -- 16 bit}}\\
\hline
\multicolumn{3}{|c|}{\cellcolor{acikGri}} & \cellcolor{turuncu!25}\textbf{AH (üst 8)}\ \textbar\ \cellcolor{yesil!25}\textbf{AL (alt 8)}\\
\hline
\end{tabular}
\end{center}

\begin{center}
\small
\begin{tabularx}{\textwidth}{@{}l l l l >{\raggedright\arraybackslash}X@{}}
\toprule
\rowcolor{anaMavi}
\color{white}\textbf{64-bit} & \color{white}\textbf{32-bit} & \color{white}\textbf{16-bit} & \color{white}\textbf{8-bit} & \color{white}\textbf{Geleneksel Amaç} \\
\midrule
\reg{RAX} & \reg{EAX} & \reg{AX} & \reg{AH}/\reg{AL} & \textbf{Accumulator} -- aritmetik, çarpma/bölme, syscall dönüş değeri \\
\rowcolor{acikGri}
\reg{RBX} & \reg{EBX} & \reg{BX} & \reg{BH}/\reg{BL} & \textbf{Base} -- bellek adreslemede taban işaretçisi \\
\reg{RCX} & \reg{ECX} & \reg{CX} & \reg{CH}/\reg{CL} & \textbf{Counter} -- döngü ve string işlemlerinde sayaç \\
\rowcolor{acikGri}
\reg{RDX} & \reg{EDX} & \reg{DX} & \reg{DH}/\reg{DL} & \textbf{Data} -- G/Ç portları, çarpma/bölmede yüksek bölüm \\
\reg{RSI} & \reg{ESI} & \reg{SI} & \reg{SIL} & \textbf{Source Index} -- string kopyalamada kaynak \\
\rowcolor{acikGri}
\reg{RDI} & \reg{EDI} & \reg{DI} & \reg{DIL} & \textbf{Destination Index} -- string kopyalamada hedef \\
\reg{RBP} & \reg{EBP} & \reg{BP} & \reg{BPL} & \textbf{Base Pointer} -- yığın çerçevesi (stack frame) tabanı \\
\rowcolor{acikGri}
\reg{RSP} & \reg{ESP} & \reg{SP} & \reg{SPL} & \textbf{Stack Pointer} -- yığının tepesini gösterir \\
\midrule
\reg{R8}--\reg{R15} & \reg{R8D}--\reg{R15D} & \reg{R8W}--\reg{R15W} & \reg{R8B}--\reg{R15B} & 64-bit modda eklenen 8 ek register (özel amacı yok) \\
\bottomrule
\end{tabularx}
\end{center}

\begin{uyarikutu}
64-bit modda bir register'ın 32-bit parçasına yazmak (örn. \reg{EAX}), üst 32
biti \textbf{otomatik olarak sıfırlar}. Ancak 16-bit (\reg{AX}) veya 8-bit
(\reg{AL}) parçasına yazmak üst bitleri \textbf{korur}. Örnek: \reg{RAX}=\hx{0xFFFFFFFFFFFFFFFF}
iken \kod{mov eax, 0} yaparsanız \reg{RAX}=\hx{0} olur; ama \kod{mov ax, 0}
yaparsanız \reg{RAX}=\hx{0xFFFFFFFFFFFF0000} kalır. Bu, çok yaygın bir hata
kaynağıdır.
\end{uyarikutu}

\subsection{Her Register'ın Detaylı İşlevi ve Örnekleri}

\subsubsection{RAX / EAX / AX -- Accumulator (Biriktirici)}
Aritmetik işlemlerin favori register'ıdır. \kod{MUL} ve \kod{DIV} komutları
örtük olarak bu register'ı kullanır. Linux'ta \kod{syscall} numarası buraya
yazılır ve dönüş değeri buradan okunur.

\begin{lstlisting}
mov eax, 5        ; EAX = 5
add eax, 3        ; EAX = 8
imul eax, 4       ; EAX = 32  (isaretli carpma)
mov al, 0xFF      ; sadece alt bayti degistir
                  ; EAX'in ust 24 biti korunur: EAX = 0x000000FF
\end{lstlisting}

\subsubsection{RBX / EBX / BX -- Base (Taban)}
Genellikle bir bellek bloğunun/dizinin başlangıç adresini tutmak için kullanılır.
Gerçek modda dolaylı adreslemede sık kullanılır.

\begin{lstlisting}
mov bx, dizi      ; BX = dizinin baslangic adresi
mov al, [bx]      ; AL = dizinin ilk elemani
inc bx            ; sonraki elemana gec
mov al, [bx]      ; AL = ikinci eleman
\end{lstlisting}

\subsubsection{RCX / ECX / CX -- Counter (Sayaç)}
\kod{LOOP} komutu her çalıştığında \reg{CX}/\reg{ECX}'i bir azaltır ve sıfır
olana kadar döngüye devam eder. \kod{REP} önekiyle string komutlarında tekrar
sayısını, \kod{SHL}/\kod{SHR} komutlarında ise kaydırma miktarını belirler.

\begin{lstlisting}
mov ecx, 5        ; 5 kez donecek
etiket:
    ; ... yapilacak is ...
    loop etiket   ; ECX--, sifir degilse etiket'e don

; SHL/SHR icin CL kullanimi:
mov al, 1
mov cl, 4
shl al, cl        ; AL = 1 << 4 = 16
\end{lstlisting}

\subsubsection{RDX / EDX / DX -- Data (Veri)}
Çarpma ve bölmede 64/128-bit sonucun \emph{yüksek} yarısını tutar. Ayrıca
G/Ç port numarası \reg{DX}'te tutulur.

\begin{ornekkutu}[32-bit çarpma ve bölme -- RDX:RAX çifti]
\begin{lstlisting}
; CARPMA: EAX * EBX -> sonuc 64-bit olarak EDX:EAX
mov eax, 0x10000000
mov ebx, 16
mul ebx           ; EDX:EAX = EAX * EBX
                  ; EDX = ust 32 bit, EAX = alt 32 bit

; BOLME: EDX:EAX / EBX -> bolum EAX, kalan EDX
mov edx, 0        ; ONEMLI: bolmeden once EDX'i temizle!
mov eax, 100
mov ebx, 7
div ebx           ; EAX = 100/7 = 14,  EDX = 100 mod 7 = 2
\end{lstlisting}
\end{ornekkutu}

\begin{uyarikutu}
\kod{DIV} komutundan önce \reg{EDX} (veya 16-bit'te \reg{DX}) mutlaka
temizlenmelidir; aksi halde işlemci temizlenmemiş yüksek biti bölmeye dahil
eder ve genellikle "bölme taşması" (divide overflow) kesmesi (\hx{INT 0}) oluşur.
İşaretli bölmede \kod{CDQ} komutu \reg{EAX}'in işaretini \reg{EDX}'e yayar.
\end{uyarikutu}

\subsubsection{RSP \& RBP -- Yığın (Stack) Yönetimi}
\reg{RSP} yığının tepesini gösterir; \kod{PUSH} ile azalır, \kod{POP} ile
artar (yığın yukarıdan aşağıya büyür). \reg{RBP} ise bir fonksiyon içindeki
yerel değişkenlere ve parametrelere sabit bir referans noktası sağlar.

\begin{ornekkutu}[Standart fonksiyon prologu ve epilogu]
\begin{lstlisting}
topla:
    push rbp          ; 1) eski taban isaretcisini sakla
    mov rbp, rsp      ; 2) yeni cerceve tabani olustur
    sub rsp, 16       ; 3) 16 bayt yerel degisken alani ayir

    mov [rbp-8], rdi  ; ilk parametreyi yerel degiskene koy
    mov rax, [rbp-8]
    add rax, rsi      ; RAX = param1 + param2

    mov rsp, rbp      ; 4) yerel alani geri al
    pop rbp           ; 5) eski tabani geri yukle
    ret               ; 6) fonksiyondan don (donus degeri RAX'te)
\end{lstlisting}
\end{ornekkutu}

% =====================================================================
\section{Adresleme Modları}
% =====================================================================

Bir komutun operandına nasıl ulaşacağını "adresleme modu" belirler. x86'nın
esnek adresleme modları, dizilere ve yapılara erişimi kolaylaştırır.

\begin{center}
\begin{tabularx}{\textwidth}{@{}l l >{\raggedright\arraybackslash}X@{}}
\toprule
\rowcolor{anaMavi}\color{white}\textbf{Mod} & \color{white}\textbf{Örnek} & \color{white}\textbf{Anlamı} \\
\midrule
Anlık (immediate) & \kod{mov ax, 5} & Sabit değer doğrudan komutun içinde \\
\rowcolor{acikGri} Register & \kod{mov ax, bx} & Değer bir register'dan alınır \\
Doğrudan (direct) & \kod{mov ax, [1234h]} & Belirli bir bellek adresi \\
\rowcolor{acikGri} Register dolaylı & \kod{mov ax, [bx]} & BX'in gösterdiği adresteki değer \\
Taban + yer değiştirme & \kod{mov ax, [bx+4]} & Yapı alanlarına erişim \\
\rowcolor{acikGri} Taban + indeks & \kod{mov ax, [bx+si]} & İki register toplamı \\
Ölçekli indeks (SIB) & \kod{mov eax, [ebx+esi*4]} & Dizi elemanına erişim (4 baytlık) \\
\bottomrule
\end{tabularx}
\end{center}

\begin{ornekkutu}[Ölçekli indeks ile bir int dizisinde gezinme]
\begin{lstlisting}
; int32 dizisi[i] elemanina erisim: adres = taban + i*4
mov ebx, dizi     ; EBX = dizinin baslangici
mov esi, 3        ; ESI = istenen indeks (i = 3)
mov eax, [ebx+esi*4]  ; EAX = dizi[3]
                      ; olcek 4 cunku her eleman 4 bayt
\end{lstlisting}
\end{ornekkutu}

\begin{bilgikutu}
Genel formül: \kod{[taban + indeks*ölçek + yer\_değiştirme]}. Ölçek yalnızca
1, 2, 4 veya 8 olabilir (byte, word, dword, qword eleman boyutlarına karşılık
gelir). Bu tek komut, C dilindeki \kod{dizi[i]} ifadesinin doğrudan
karşılığıdır.
\end{bilgikutu}

\begin{ipucukutu}
\kod{LEA} (Load Effective Address) komutu, adresi \emph{hesaplar} ama belleğe
gitmez; sadece hesaplanan adresi register'a yazar. Bu yüzden hızlı aritmetik
için de kullanılır: \kod{lea eax, [ebx+ebx*4]} aslında \kod{EAX = EBX*5}
yapar, üstelik bayrakları değiştirmeden!
\end{ipucukutu}

% =====================================================================
\section{Segment Register'ları}
% =====================================================================

Segment register'ları 16-bit'tir ve gerçek modda belleği 64\,KB'lık bloklara
(segment) böler. Fiziksel adres \kod{segment $\times$ 16 + offset} formülüyle
hesaplanır.

\begin{center}
\begin{tabularx}{\textwidth}{@{}>{\bfseries\color{anaMavi}}l l >{\raggedright\arraybackslash}X@{}}
\toprule
\rowcolor{acikMavi}\color{anaMavi}\textbf{Register} & \color{anaMavi}\textbf{Açılım} & \color{anaMavi}\textbf{İşlev} \\
\midrule
CS & Code Segment & Çalışan kodun bulunduğu segment (IP ile eşleşir) \\
\rowcolor{acikGri} DS & Data Segment & Varsayılan veri segmenti \\
SS & Stack Segment & Yığının bulunduğu segment (SP/BP ile eşleşir) \\
\rowcolor{acikGri} ES & Extra Segment & Ek veri segmenti (string hedefi için DI ile) \\
FS & -- & Ek segment (modern OS'lerde thread-local storage) \\
\rowcolor{acikGri} GS & -- & Ek segment (64-bit'te çekirdek veri yapıları) \\
\bottomrule
\end{tabularx}
\end{center}

\begin{ornekkutu}[Fiziksel Adres Hesabı ve segment kurulumu]
\reg{DS} = \hx{0x1000}, offset = \hx{0x0234} ise\\
Fiziksel adres = \hx{0x1000} $\times$ 16 + \hx{0x0234} = \hx{0x10000} + \hx{0x0234} = \hx{0x10234}
\begin{lstlisting}
; Onyukleyicide segmentleri guvenli bir sekilde ayarlama
xor ax, ax        ; AX = 0
mov ds, ax        ; DS = 0  (segmentlere dogrudan sabit yazilamaz,
mov es, ax        ; ES = 0   once bir register'a koymak gerekir)
mov ss, ax        ; SS = 0
mov sp, 0x7C00    ; yigin onyukleyicinin hemen altinda
\end{lstlisting}
\end{ornekkutu}

\begin{bilgikutu}
Korumalı ve uzun modda segment register'ları artık bellek bölme için değil,
\textbf{segment seçici (selector)} olarak kullanılır ve GDT/LDT tablolarındaki
tanımlayıcıları işaret eder. 64-bit'te \reg{FS} ve \reg{GS} dışındaki
segmentlerin tabanı genellikle 0 kabul edilir (flat / düz model).
\end{bilgikutu}

% =====================================================================
\section{İşaretçi ve İndeks Register'ları}
% =====================================================================

\begin{center}
\begin{tabularx}{\textwidth}{@{}>{\bfseries\color{anaMavi}}l l >{\raggedright\arraybackslash}X@{}}
\toprule
\rowcolor{acikMavi}\color{anaMavi}\textbf{Register} & \color{anaMavi}\textbf{Ad} & \color{anaMavi}\textbf{İşlev} \\
\midrule
IP / EIP / RIP & Instruction Pointer & Bir sonraki çalıştırılacak komutun adresi \\
\rowcolor{acikGri} SP / ESP / RSP & Stack Pointer & Yığının tepesi \\
BP / EBP / RBP & Base Pointer & Yığın çerçevesi tabanı \\
\rowcolor{acikGri} SI / ESI / RSI & Source Index & String/dizi kaynağı \\
DI / EDI / RDI & Destination Index & String/dizi hedefi \\
\bottomrule
\end{tabularx}
\end{center}

\subsection{SI ve DI ile String İşlemleri}

\reg{SI} ve \reg{DI}, \kod{MOVSB}, \kod{LODSB}, \kod{STOSB} gibi string
komutlarında otomatik olarak kullanılır. Bu komutlar \reg{DF} (Direction Flag)
bayrağına göre indeksleri otomatik artırır/azaltır.

\begin{ornekkutu}[Bir bellek bloğunu kopyalama (memcpy)]
\begin{lstlisting}
cld               ; DF=0 -> SI ve DI artacak (ileri yon)
mov si, kaynak    ; DS:SI = kaynak adres
mov di, hedef     ; ES:DI = hedef adres
mov cx, 100       ; 100 bayt kopyalanacak
rep movsb         ; CX kez: [ES:DI] <- [DS:SI], SI++, DI++, CX--
\end{lstlisting}
Tek satırlık \kod{rep movsb} ile 100 baytlık kopyalama tamamlanır. Bu, C'deki
\kod{memcpy(hedef, kaynak, 100)} ile eşdeğerdir.
\end{ornekkutu}

\begin{ornekkutu}[Bir alanı sabit değerle doldurma (memset)]
\begin{lstlisting}
cld
mov di, tampon    ; ES:DI = doldurulacak alan
mov al, 0         ; doldurulacak deger (0)
mov cx, 256       ; 256 bayt
rep stosb         ; CX kez: [ES:DI] <- AL, DI++, CX--
\end{lstlisting}
\end{ornekkutu}

\begin{uyarikutu}
\reg{RIP} (Instruction Pointer) doğrudan \kod{mov} ile değiştirilemez.
Yalnızca \kod{jmp}, \kod{call}, \kod{ret} gibi akış komutlarıyla dolaylı olarak
değişir. 64-bit modda \reg{RIP}-göreli adresleme (RIP-relative), konumdan
bağımsız kod (PIC) yazmak için kullanılır.
\end{uyarikutu}

% =====================================================================
\section{Bayrak Register'ı (FLAGS / EFLAGS / RFLAGS)}
% =====================================================================

Bayrak register'ı, işlemcinin durum bilgilerini tutan tekil bitlerden oluşur.
Aritmetik/mantık komutları bu bayrakları günceller; koşullu atlamalar
(\kod{JE}, \kod{JG}, vb.) bunları okuyarak karar verir.

\begin{center}
\small
\begin{tabularx}{\textwidth}{@{}c l l >{\raggedright\arraybackslash}X@{}}
\toprule
\rowcolor{anaMavi}
\color{white}\textbf{Bit} & \color{white}\textbf{Kısa} & \color{white}\textbf{Ad} & \color{white}\textbf{Anlamı} \\
\midrule
0 & CF & Carry Flag & Elde/ödünç oluştuğunda 1 (işaretsiz taşma) \\
\rowcolor{acikGri} 2 & PF & Parity Flag & Sonucun düşük baytında çift sayıda 1 biti varsa set \\
4 & AF & Auxiliary Carry & BCD aritmetiğinde yarım-bayt eldesi \\
\rowcolor{acikGri} 6 & ZF & Zero Flag & Sonuç sıfır ise 1 \\
7 & SF & Sign Flag & Sonucun en yüksek biti (negatiflik) \\
\rowcolor{acikGri} 8 & TF & Trap Flag & Tek adım (single-step) hata ayıklama modu \\
9 & IF & Interrupt Enable & Donanım kesmeleri açık (1) / kapalı (0) \\
\rowcolor{acikGri} 10 & DF & Direction Flag & String işlemleri yönü: 0 artan, 1 azalan \\
11 & OF & Overflow Flag & İşaretli aritmetikte taşma oluştuğunda 1 \\
\bottomrule
\end{tabularx}
\end{center}

\subsection{Koşullu Atlama Komutları}

Karşılaştırmadan (\kod{CMP} veya \kod{TEST}) sonra kullanılan atlama komutları.
İşaretli ve işaretsiz sayılar için ayrı komutlar olduğuna dikkat edin.

\begin{center}
\small
\begin{tabularx}{\textwidth}{@{}l l l l@{}}
\toprule
\rowcolor{anaMavi}
\color{white}\textbf{Komut} & \color{white}\textbf{Koşul} & \color{white}\textbf{Bayrak} & \color{white}\textbf{Tür} \\
\midrule
JE / JZ & Eşit / Sıfır & ZF=1 & -- \\
\rowcolor{acikGri} JNE / JNZ & Eşit değil & ZF=0 & -- \\
JG / JNLE & Büyük & ZF=0 ve SF=OF & işaretli \\
\rowcolor{acikGri} JL / JNGE & Küçük & SF$\neq$OF & işaretli \\
JGE / JNL & Büyük eşit & SF=OF & işaretli \\
\rowcolor{acikGri} JLE / JNG & Küçük eşit & ZF=1 veya SF$\neq$OF & işaretli \\
JA / JNBE & Büyük (above) & CF=0 ve ZF=0 & işaretsiz \\
\rowcolor{acikGri} JB / JC & Küçük (below) & CF=1 & işaretsiz \\
JAE / JNC & Büyük eşit & CF=0 & işaretsiz \\
\rowcolor{acikGri} JBE / JNA & Küçük eşit & CF=1 veya ZF=1 & işaretsiz \\
\bottomrule
\end{tabularx}
\end{center}

\begin{ornekkutu}[Karşılaştırma ve dallanma -- basit if/else]
\begin{lstlisting}
mov eax, [sayi]
cmp eax, 10       ; EAX - 10 hesapla, bayraklari ayarla
jg  buyuk         ; isaretli "buyuk mu?" -> EAX > 10 ise atla
    ; buraya gelindiyse EAX <= 10
    mov ebx, 0
    jmp bitti
buyuk:
    mov ebx, 1    ; EAX > 10 durumu
bitti:
\end{lstlisting}
\end{ornekkutu}

\begin{ipucukutu}
\kod{TEST eax, eax} bir sayının sıfır olup olmadığını kontrol etmenin en hızlı
yoludur (\kod{cmp eax, 0}'dan daha verimli). \kod{CLI}/\kod{STI} kesmeleri
kapatır/açar; \kod{CLD}/\kod{STD} ise string yönünü kontrol eder.
\end{ipucukutu}

% =====================================================================
\section{SIMD ve FPU Register'ları}
% =====================================================================

Modern işlemciler, tek komutla birden fazla veri üzerinde işlem yapan
(SIMD -- Single Instruction Multiple Data) özel register setlerine sahiptir.
Bunlar grafik, ses, bilimsel hesaplama ve şifrelemede yoğun kullanılır.

\begin{center}
\begin{tabularx}{\textwidth}{@{}l l >{\raggedright\arraybackslash}X@{}}
\toprule
\rowcolor{anaMavi}\color{white}\textbf{Set} & \color{white}\textbf{Register'lar} & \color{white}\textbf{Açıklama} \\
\midrule
x87 FPU & ST0--ST7 (80-bit) & Kayan noktalı (floating point) yığın tabanlı hesaplama \\
\rowcolor{acikGri} MMX & MM0--MM7 (64-bit) & Tamsayı SIMD (FPU register'larıyla örtüşür) \\
SSE & XMM0--XMM15 (128-bit) & 4$\times$float veya 2$\times$double paralel işlem \\
\rowcolor{acikGri} AVX & YMM0--YMM15 (256-bit) & 8$\times$float paralel işlem \\
AVX-512 & ZMM0--ZMM31 (512-bit) & 16$\times$float, maske register'ları (K0--K7) \\
\bottomrule
\end{tabularx}
\end{center}

\begin{bilgikutu}
XMM, YMM ve ZMM register'ları iç içedir: XMM0, YMM0'ın alt 128 biti; YMM0 da
ZMM0'ın alt 256 bitidir. Tıpkı AX'in EAX'in içinde olması gibi. Örneğin bir
\kod{addps xmm0, xmm1} komutu, iki register'daki 4'er float değeri
\emph{aynı anda} toplar.
\end{bilgikutu}

\begin{ornekkutu}[SSE ile 4 float'ı aynı anda toplama]
\begin{lstlisting}
; a[4] ve b[4] float dizilerini topla -> sonuc[4]
movups xmm0, [a]  ; XMM0 = {a0, a1, a2, a3}  (128-bit yukle)
movups xmm1, [b]  ; XMM1 = {b0, b1, b2, b3}
addps  xmm0, xmm1 ; XMM0 = {a0+b0, a1+b1, a2+b2, a3+b3} tek komutta!
movups [sonuc], xmm0
\end{lstlisting}
Skaler kodda 4 ayrı toplama gerekirken SSE bunu tek komutta yapar -- bu
\emph{paralellik} SIMD'in gücüdür.
\end{ornekkutu}

% =====================================================================
\section{Kontrol ve Özel Register'lar}
% =====================================================================

Bunlar işletim sistemi seviyesinde, işlemcinin çalışma modunu ve bellek
yönetimini kontrol eden ayrıcalıklı register'lardır (yalnızca ring 0).

\begin{center}
\begin{tabularx}{\textwidth}{@{}>{\bfseries\color{anaMavi}}l >{\raggedright\arraybackslash}X@{}}
\toprule
\rowcolor{acikMavi}\color{anaMavi}\textbf{Register} & \color{anaMavi}\textbf{İşlev} \\
\midrule
CR0 & Korumalı moda geçiş (bit 0 = PE), sayfalama (bit 31 = PG), önbellek kontrolü \\
\rowcolor{acikGri} CR2 & Sayfa hatası (page fault) oluştuğunda erişilmeye çalışılan adres \\
CR3 & Sayfa dizini taban adresi (PDBR) -- sanal bellek haritası \\
\rowcolor{acikGri} CR4 & PAE, PSE, SSE, PCID gibi gelişmiş özelliklerin etkinleştirilmesi \\
CR8 & (64-bit) Görev önceliği register'ı (TPR) \\
\rowcolor{acikGri} GDTR / IDTR & Global / Interrupt tanımlayıcı tablo adres ve boyutları \\
DR0--DR7 & Donanım kesme noktaları (hardware breakpoints) için hata ayıklama register'ları \\
MSR & Model'e özel register'lar (\kod{rdmsr}/\kod{wrmsr} ile erişilir) \\
\bottomrule
\end{tabularx}
\end{center}

% =====================================================================
\section{Yaygın Komut Hızlı Referansı}
% =====================================================================

\begin{center}
\small
\begin{tabularx}{\textwidth}{@{}l >{\raggedright\arraybackslash}X l >{\raggedright\arraybackslash}X@{}}
\toprule
\rowcolor{anaMavi}\color{white}\textbf{Komut} & \color{white}\textbf{İşlev} & \color{white}\textbf{Komut} & \color{white}\textbf{İşlev} \\
\midrule
MOV & Değer taşı & CMP & Karşılaştır (çıkar, bayrak ayarla) \\
\rowcolor{acikGri} ADD/SUB & Topla / Çıkar & TEST & Bit AND (bayrak ayarla) \\
INC/DEC & 1 artır / azalt & AND/OR/XOR & Bit işlemleri \\
\rowcolor{acikGri} MUL/IMUL & İşaretsiz / işaretli çarpma & NOT/NEG & Bit tümleyeni / negatif \\
DIV/IDIV & İşaretsiz / işaretli bölme & SHL/SHR & Sola / sağa kaydırma \\
\rowcolor{acikGri} PUSH/POP & Yığına at / çek & ROL/ROR & Döndürme (rotate) \\
CALL/RET & Fonksiyon çağır / dön & LEA & Etkin adres yükle \\
\rowcolor{acikGri} JMP & Koşulsuz atlama & MOVZX/MOVSX & Sıfır / işaret genişletmeli taşıma \\
INT & Yazılım kesmesi & NOP & Boş işlem \\
\rowcolor{acikGri} IN/OUT & G/Ç portundan oku/yaz & HLT & İşlemciyi durdur \\
\bottomrule
\end{tabularx}
\end{center}

\begin{ornekkutu}[Bit işlemleriyle yaygın hileler]
\begin{lstlisting}
xor eax, eax      ; EAX = 0  (en hizli sifirlama yontemi)
test eax, eax     ; EAX sifir mi? -> ZF ayarlanir
and al, 0x0F      ; sadece alt 4 biti tut (maskeleme)
or  al, 0x80      ; en yuksek biti 1 yap
shl eax, 1        ; 2 ile carp (sola 1 kaydir)
shr eax, 2        ; 4 e bol (saga 2 kaydir, isaretsiz)
\end{lstlisting}
\end{ornekkutu}

% =====================================================================
\section{BIOS Interrupt (Kesme) Sistemi}
% =====================================================================

BIOS (Basic Input/Output System), anakarttaki bir çip üzerinde bulunan ve
donanıma en alt seviyede erişim sağlayan bir yazılımdır. \textbf{Kesmeler}
(interrupt), donanım hizmetlerine erişmek için standartlaştırılmış giriş
noktalarıdır.

\subsection{Kesme Nasıl Çalışır? -- Kesme Vektör Tablosu}

Bir yazılım kesmesi \kod{INT n} komutuyla tetiklenir. İşlemci, belleğin en
başında (\hx{0x0000}) bulunan \textbf{Kesme Vektör Tablosu}'na (IVT) bakar.
Bu tabloda her girdi 4 bayttır (2 bayt offset + 2 bayt segment). İşlemci
\kod{n$\times$4} adresindeki segment:offset değerine atlar.

\begin{center}
\begin{tabular}{|c|c|c|c|}
\hline
\rowcolor{anaMavi}\multicolumn{4}{|c|}{\color{white}\textbf{Kesme Vektör Tablosu (IVT) -- adres 0x0000}}\\
\hline
\cellcolor{acikMavi}\textbf{Girdi} & \cellcolor{acikMavi}\textbf{Adres} & \cellcolor{acikMavi}\textbf{İçerik} & \cellcolor{acikMavi}\textbf{Kesme} \\
\hline
0 & 0x0000 & offset:segment & INT 0 (bölme hatası) \\
\hline
1 & 0x0004 & offset:segment & INT 1 (tek adım) \\
\hline
\dots & \dots & \dots & \dots \\
\hline
16 & 0x0040 & offset:segment & INT 10h (video) \\
\hline
\end{tabular}
\end{center}

\begin{bilgikutu}
\textbf{Genel Kullanım Deseni:} Kesme numarası hangi \emph{cihazı/hizmeti}
seçer. \reg{AH} register'ı ise o hizmet içindeki hangi \emph{alt fonksiyonun}
çağrılacağını belirler. Diğer register'lar (\reg{AL}, \reg{BX}, \reg{CX},
\reg{DX}) parametreleri taşır. Sonuç genellikle register'larda ve/veya
\reg{CF} (Carry) bayrağında döner: \kod{CF=1} çoğunlukla "hata oluştu"
anlamına gelir.
\end{bilgikutu}

\subsection{Önemli BIOS Kesmelerine Genel Bakış}

\begin{center}
\begin{tabularx}{\textwidth}{@{}c l >{\raggedright\arraybackslash}X@{}}
\toprule
\rowcolor{anaMavi}
\color{white}\textbf{INT} & \color{white}\textbf{Hizmet} & \color{white}\textbf{Açıklama} \\
\midrule
\hx{10h} & Video Servisleri & Ekrana yazı yazma, piksel çizme, mod ayarı, imleç \\
\rowcolor{acikGri} \hx{11h} & Ekipman Listesi & Takılı donanımların bit maskesi \\
\hx{12h} & Bellek Boyutu & Kullanılabilir taban bellek (KB) \\
\rowcolor{acikGri} \hx{13h} & Disk Servisleri & Sektör okuma/yazma, disk resetleme, LBA \\
\hx{14h} & Seri Port & RS-232 haberleşme \\
\rowcolor{acikGri} \hx{15h} & Sistem Servisleri & Genişletilmiş bellek haritası (E820), A20, APM \\
\hx{16h} & Klavye Servisleri & Tuş okuma, klavye durumu \\
\rowcolor{acikGri} \hx{17h} & Yazıcı Servisleri & Paralel port yazıcı \\
\hx{19h} & Bootstrap & Yeniden önyükleme \\
\rowcolor{acikGri} \hx{1Ah} & Saat Servisleri & Sistem saati, RTC okuma \\
\bottomrule
\end{tabularx}
\end{center}

% =====================================================================
\section{INT 10h -- Video Servisleri (Detaylı)}
% =====================================================================

Ekranla ilgili tüm işlemler bu kesme üzerinden yapılır.

\subsection{AH=0Eh -- Teletype Modunda Karakter Yazma}

Bu fonksiyon, imleci otomatik ilerleterek tek bir karakteri ekrana basar.
Önyükleyicilerde "Hello World" yazdırmanın standart yoludur.

\begin{center}
\begin{tabularx}{\textwidth}{@{}l >{\raggedright\arraybackslash}X@{}}
\toprule
\rowcolor{acikMavi}\color{anaMavi}\textbf{Giriş} & \color{anaMavi}\textbf{Anlam} \\
\midrule
\reg{AH} = \hx{0Eh} & Fonksiyon: teletype yazma \\
\reg{AL} = karakter & Yazdırılacak ASCII karakter kodu \\
\reg{BH} = sayfa & Görüntü sayfası (genellikle 0) \\
\reg{BL} = renk & Grafik modunda ön plan rengi \\
\bottomrule
\end{tabularx}
\end{center}

\begin{ornekkutu}[Null ile biten metin yazan yeniden kullanılabilir rutin]
\begin{lstlisting}
yaz_metin:            ; giris: SI = metnin adresi
    push ax
    mov ah, 0x0E      ; teletype fonksiyonu
.dongu:
    lodsb             ; DS:SI'daki bayti AL'e yukle, SI++
    test al, al       ; AL == 0 mi? (metin sonu)
    jz  .bitti
    int 0x10          ; karakteri ekrana yaz
    jmp .dongu
.bitti:
    pop ax
    ret

; kullanimi:
mov si, mesaj
call yaz_metin
mesaj: db "Merhaba, Assembly Dunyasi!", 0x0D, 0x0A, 0
\end{lstlisting}
\end{ornekkutu}

\begin{notkutu}[Satır sonu karakterleri]
\hx{0x0D} (Carriage Return -- satır başı) ve \hx{0x0A} (Line Feed -- alt satır)
birlikte kullanılır. Sadece \hx{0x0A} imleci alt satıra indirir ama başa
getirmez; bu yüzden ikisi birlikte gerekir.
\end{notkutu}

\subsection{AH=00h -- Video Modu Ayarlama}

\begin{center}
\begin{tabularx}{\textwidth}{@{}l l >{\raggedright\arraybackslash}X@{}}
\toprule
\rowcolor{acikMavi}\color{anaMavi}\textbf{AL} & \color{anaMavi}\textbf{Mod} & \color{anaMavi}\textbf{Açıklama} \\
\midrule
\hx{03h} & 80$\times$25 metin & 16 renkli standart metin modu \\
\rowcolor{acikGri} \hx{12h} & 640$\times$480 & 16 renkli grafik (VGA) \\
\hx{13h} & 320$\times$200 & 256 renkli grafik (oyunların klasiği) \\
\bottomrule
\end{tabularx}
\end{center}

\begin{lstlisting}
mov ah, 0x00      ; video modu ayarla
mov al, 0x03      ; 80x25 metin modu (ekrani da temizler)
int 0x10
\end{lstlisting}

\subsection{AH=02h -- İmleç Konumu Ayarlama}

\begin{center}
\begin{tabularx}{\textwidth}{@{}l >{\raggedright\arraybackslash}X@{}}
\toprule
\rowcolor{acikMavi}\color{anaMavi}\textbf{Giriş} & \color{anaMavi}\textbf{Anlam} \\
\midrule
\reg{AH}=\hx{02h}, \reg{BH}=sayfa & İmleç konumu ayarla \\
\reg{DH} = satır & 0--24 arası (üstten) \\
\reg{DL} = sütun & 0--79 arası (soldan) \\
\bottomrule
\end{tabularx}
\end{center}

\begin{ornekkutu}[İmleci 10. satır, 20. sütuna taşıma]
\begin{lstlisting}
mov ah, 0x02
mov bh, 0         ; sayfa 0
mov dh, 10        ; satir 10
mov dl, 20        ; sutun 20
int 0x10
\end{lstlisting}
\end{ornekkutu}

\subsection{AH=09h -- Karakteri Renkli ve Tekrarlı Yazma}

İmleci ilerletmeden, belirtilen karakteri belirtilen renkte ve sayıda yazar.
Renkli arayüzler çizmek için idealdir.

\begin{center}
\begin{tabularx}{\textwidth}{@{}l >{\raggedright\arraybackslash}X@{}}
\toprule
\rowcolor{acikMavi}\color{anaMavi}\textbf{Giriş} & \color{anaMavi}\textbf{Anlam} \\
\midrule
\reg{AH}=\hx{09h}, \reg{AL}=karakter & Yazılacak karakter \\
\reg{BH}=sayfa, \reg{BL}=öznitelik & Renk: alt 4 bit ön plan, üst 4 bit arka plan \\
\reg{CX} = tekrar sayısı & Kaç kez yazılacak \\
\bottomrule
\end{tabularx}
\end{center}

\begin{ornekkutu}[Beyaz zemin üzerine kırmızı 10 yıldız]
\begin{lstlisting}
mov ah, 0x09
mov al, '*'       ; yildiz karakteri
mov bh, 0         ; sayfa 0
mov bl, 0xF4      ; F=beyaz arka plan, 4=kirmizi on plan
mov cx, 10        ; 10 kez
int 0x10
\end{lstlisting}
\end{ornekkutu}

\begin{bilgikutu}
\textbf{Renk kodları (4-bit):}
0=Siyah, 1=Mavi, 2=Yeşil, 3=Cyan, 4=Kırmızı, 5=Magenta, 6=Kahve, 7=Açık gri,
8=Koyu gri, 9=Açık mavi, A=Açık yeşil, B=Açık cyan, C=Açık kırmızı,
D=Açık magenta, E=Sarı, F=Beyaz.
\end{bilgikutu}

\subsection{AH=06h -- Ekranı Kaydırma / Temizleme}

Bir pencere alanını yukarı kaydırır. \reg{AL}=0 verildiğinde tüm alanı
temizler -- ekranı belirli renkle temizlemenin yaygın yolu.

\begin{ornekkutu}[Tüm ekranı mavi zeminle temizleme]
\begin{lstlisting}
mov ah, 0x06      ; yukari kaydir / temizle
mov al, 0         ; 0 = tum pencereyi temizle
mov bh, 0x1F      ; renk: 1=mavi arka, F=beyaz on plan
mov cx, 0x0000    ; sol-ust kose (satir 0, sutun 0)
mov dx, 0x184F    ; sag-alt kose (satir 24=0x18, sutun 79=0x4F)
int 0x10
\end{lstlisting}
\end{ornekkutu}

\subsection{AH=0Ch -- Piksel Çizme (Grafik Modu)}

\begin{center}
\begin{tabularx}{\textwidth}{@{}l >{\raggedright\arraybackslash}X@{}}
\toprule
\rowcolor{acikMavi}\color{anaMavi}\textbf{Giriş} & \color{anaMavi}\textbf{Anlam} \\
\midrule
\reg{AH}=\hx{0Ch}, \reg{AL}=renk & Renk numarası (0--255) \\
\reg{CX} = X koordinatı & Sütun \\
\reg{DX} = Y koordinatı & Satır \\
\bottomrule
\end{tabularx}
\end{center}

\begin{ornekkutu}[Grafik modda yatay bir kırmızı çizgi çizme]
\begin{lstlisting}
mov ah, 0x00      ; once grafik moduna gec
mov al, 0x13      ; 320x200, 256 renk
int 0x10

mov ah, 0x0C      ; piksel ciz
mov al, 4         ; renk: kirmizi
mov dx, 100       ; Y = 100 (sabit satir)
mov cx, 0         ; X = 0 den basla
.ciz:
    int 0x10      ; (CX, DX) noktasina piksel koy
    inc cx        ; X++
    cmp cx, 320   ; ekran genisligine ulastik mi?
    jl  .ciz      ; hayirsa devam
\end{lstlisting}
\end{ornekkutu}

% =====================================================================
\section{INT 16h -- Klavye Servisleri (Detaylı)}
% =====================================================================

\subsection{AH=00h -- Tuş Bekle ve Oku (Bloke Eder)}

Bir tuşa basılana kadar programı bekletir. \reg{AL}'de ASCII kodu, \reg{AH}'de
tarama kodu (scan code) döner.

\begin{ornekkutu}[Kullanıcıdan tuş bekleyip ekrana yankılama (echo terminali)]
\begin{lstlisting}
bekle_tus:
    mov ah, 0x00  ; tus bekle ve oku
    int 0x16      ; AL = basilan tusun ASCII kodu

    cmp al, 0x0D  ; Enter'a mi basildi?
    je  yeni_satir

    mov ah, 0x0E  ; teletype yazma
    int 0x10      ; okunan tusu ekrana yaz
    jmp bekle_tus

yeni_satir:
    mov ah, 0x0E
    mov al, 0x0D
    int 0x10
    mov al, 0x0A
    int 0x10
    jmp bekle_tus
\end{lstlisting}
\end{ornekkutu}

\subsection{AH=01h -- Tuş Var mı? (Bloke Etmez)}

Tampon boşsa \reg{ZF}=1 döner (tuş yok); doluysa \reg{ZF}=0 olur ve \reg{AL}'de
tuş bilgisi bulunur ama \emph{tampondan silinmez}. Oyun döngülerinde giriş
kontrolü için idealdir.

\begin{ornekkutu}[Oyun döngüsünde bloke etmeden tuş kontrolü]
\begin{lstlisting}
oyun_dongusu:
    ; ... oyun mantigini guncelle, ekrani ciz ...

    mov ah, 0x01  ; tampon dolu mu?
    int 0x16
    jz  oyun_dongusu  ; ZF=1 -> tus yok, donguye devam

    ; buraya gelindiyse bir tus var, simdi tampondan al
    mov ah, 0x00
    int 0x16      ; AL = tus
    cmp al, 'q'   ; cikis tusu mu?
    jne oyun_dongusu
    ; q'ya basildi -> oyundan cik
\end{lstlisting}
\end{ornekkutu}

% =====================================================================
\section{INT 13h -- Disk Servisleri (Detaylı)}
% =====================================================================

Diskten sektör okuma, önyükleyicinin çekirdeği (kernel) belleğe yüklemesi için
kritiktir. Adresleme CHS (Cylinder-Head-Sector) veya modern LBA yöntemiyle
yapılır.

\subsection{AH=00h -- Disk Sistemini Resetleme}

Okuma hatalarından sonra sürücüyü sıfırlamak için kullanılır. Yeniden deneme
mantığının parçasıdır.

\begin{lstlisting}
mov ah, 0x00      ; disk sistemini resetle
mov dl, 0x80      ; ilk sabit disk
int 0x13
\end{lstlisting}

\subsection{AH=02h -- Sektör Okuma (CHS)}

\begin{center}
\begin{tabularx}{\textwidth}{@{}l >{\raggedright\arraybackslash}X@{}}
\toprule
\rowcolor{acikMavi}\color{anaMavi}\textbf{Giriş} & \color{anaMavi}\textbf{Anlam} \\
\midrule
\reg{AH}=\hx{02h} & Sektör oku fonksiyonu \\
\reg{AL} & Okunacak sektör sayısı \\
\reg{CH} & Silindir (cylinder) numarası \\
\reg{CL} & Sektör numarası (\textbf{1'den başlar!}) \\
\reg{DH} & Kafa (head) numarası \\
\reg{DL} & Sürücü (\hx{00h}=disket, \hx{80h}=ilk HDD) \\
\reg{ES:BX} & Verinin yükleneceği bellek adresi \\
\bottomrule
\end{tabularx}
\end{center}

\begin{center}
\begin{tabularx}{\textwidth}{@{}l >{\raggedright\arraybackslash}X@{}}
\toprule
\rowcolor{acikMavi}\color{anaMavi}\textbf{Çıkış} & \color{anaMavi}\textbf{Anlam} \\
\midrule
\reg{CF}=0 & Başarılı, \reg{AL}=okunan sektör sayısı \\
\reg{CF}=1 & Hata, \reg{AH}=hata kodu \\
\bottomrule
\end{tabularx}
\end{center}

\begin{ornekkutu}[Yeniden deneme mantığıyla sağlam sektör okuma]
\begin{lstlisting}
disk_oku:             ; ES:BX hedef, sektor bilgileri ayarli varsayilir
    mov di, 3         ; en fazla 3 kez dene
.dene:
    mov ah, 0x02      ; sektor oku
    mov al, 1         ; 1 sektor
    mov ch, 0         ; silindir 0
    mov cl, 2         ; sektor 2 (1. sektor onyukleyicidir)
    mov dh, 0         ; kafa 0
    mov dl, 0x80      ; ilk sabit disk
    int 0x13
    jnc .basarili     ; CF=0 -> okuma tamam

    ; hata olustu: diski resetle ve tekrar dene
    xor ah, ah        ; AH=0 -> reset fonksiyonu
    int 0x13
    dec di
    jnz .dene         ; deneme hakki kaldiysa tekrar
    jmp disk_hatasi   ; 3 deneme de basarisiz

.basarili:
    ret
\end{lstlisting}
\end{ornekkutu}

\begin{uyarikutu}
Sektör numarası (\reg{CL}) \textbf{1'den başlar}, 0'dan değil! İlk sektör (LBA 0)
CHS'de silindir 0, kafa 0, sektör \textbf{1}'dir. Ayrıca gerçek diskler ara
sıra okuma hatası verir; bu yüzden profesyonel önyükleyiciler daima yeniden
deneme (retry) döngüsü kullanır.
\end{uyarikutu}

\subsection{AH=42h -- Genişletilmiş Okuma (LBA)}

Modern büyük diskler için CHS yetersizdir. LBA (Logical Block Addressing)
sektörleri 0'dan başlayan tek bir numarayla adresler. Parametreler bir
"Disk Adres Paketi" (DAP) yapısı üzerinden geçer.

\begin{ornekkutu}[LBA ile sektör okuma ve DAP yapısı]
\begin{lstlisting}
    mov ah, 0x42      ; genisletilmis okuma
    mov dl, 0x80      ; disk
    mov si, dap       ; DS:SI = adres paketi
    int 0x13
    jc  disk_hatasi

; Disk Adres Paketi (16 bayt)
dap:
    db 0x10           ; paket boyutu (16)
    db 0              ; ayrilmis (0)
    dw 1              ; okunacak sektor sayisi
    dw 0x8000         ; hedef offset
    dw 0x0000         ; hedef segment
    dq 1              ; baslangic LBA numarasi (64-bit)
\end{lstlisting}
\end{ornekkutu}

% =====================================================================
\section{INT 15h, 1Ah, 12h -- Diğer Önemli Servisler}
% =====================================================================

\subsection{INT 15h, EAX=E820h -- Bellek Haritası}

Modern işletim sistemlerinin RAM haritasını öğrenmesinin standart yoludur.
Her çağrı, bir bellek bölgesini tanımlayan 20/24 baytlık bir yapı döndürür ve
\reg{EBX}'i bir sonraki bölge için günceller. \reg{EBX}=0 olunca liste biter.

\begin{ornekkutu}[E820 bellek haritası okuma döngüsü]
\begin{lstlisting}
bellek_haritasi:
    mov di, tampon    ; ES:DI = sonucun yazilacagi yer
    xor ebx, ebx      ; ilk cagri icin EBX = 0
    mov edx, 0x534D4150  ; 'SMAP' imzasi (sabit)
.dongu:
    mov eax, 0xE820   ; her cagrida yeniden ayarla
    mov ecx, 24       ; bir girdinin boyutu (bayt)
    int 0x15
    jc  .bitti        ; CF=1 -> hata veya liste sonu

    add di, 24        ; sonraki girdi icin ilerle
    test ebx, ebx     ; EBX=0 -> son girdiydi
    jnz .dongu
.bitti:
    ; tampon artik bellek bolgelerinin listesini icerir
\end{lstlisting}
\end{ornekkutu}

\begin{bilgikutu}
Her E820 girdisi şunları içerir: 8 bayt taban adres, 8 bayt uzunluk, 4 bayt
tür (1=kullanılabilir RAM, 2=rezerve, 3=ACPI, 4=NVS). İşletim sistemi bu
haritayı kullanarak hangi bellek bölgelerini güvenle kullanabileceğini belirler.
\end{bilgikutu}

\subsection{INT 15h, AX=E801h -- Genişletilmiş Bellek Boyutu}

E820'den daha basit bir alternatiftir. 15\,MB üstündeki belleği iki register
çiftinde döndürür: 1--16\,MB arası ve 16\,MB üstü.

\begin{center}
\begin{tabularx}{\textwidth}{@{}l >{\raggedright\arraybackslash}X@{}}
\toprule
\rowcolor{acikMavi}\color{anaMavi}\textbf{Çıkış} & \color{anaMavi}\textbf{Anlam} \\
\midrule
\reg{AX} / \reg{CX} & 1--16\,MB arası bellek (1\,KB birimlerle) \\
\reg{BX} / \reg{DX} & 16\,MB üstü bellek (64\,KB birimlerle) \\
\bottomrule
\end{tabularx}
\end{center}

\begin{ornekkutu}[E801 ile toplam RAM'i hesaplama]
\begin{lstlisting}
    xor cx, cx        ; CX ve BX'i temizle (bazi BIOS'lar AX/BX,
    xor bx, bx        ;   digerleri CX/DX kullanir)
    mov ax, 0xE801
    int 0x15
    jc  .hata
    ; Bazi BIOS'lar sonucu AX:BX, bazilari CX:DX doner.
    ; CX=0 ise AX/BX'i kullan:
    test cx, cx
    jnz .cxdx
    mov [dusuk_mem], ax   ; 1-16MB (KB cinsinden)
    mov [yuksek_mem], bx  ; 16MB ustu (64KB cinsinden)
    jmp .son
.cxdx:
    mov [dusuk_mem], cx
    mov [yuksek_mem], dx
.son:
.hata:
\end{lstlisting}
\end{ornekkutu}

\subsection{INT 15h, AH=88h -- Genişletilmiş Bellek (Basit)}

En eski ve en basit yöntem. \reg{AX}'te 1\,MB üstündeki bellek miktarını KB
cinsinden döndürür. Yalnızca 64\,MB'a kadar ölçeklenir, bu yüzden modern
sistemlerde E820/E801 tercih edilir.

\begin{lstlisting}
mov ah, 0x88
int 0x15          ; AX = 1MB ustundeki bellek (KB), en fazla ~64MB
\end{lstlisting}

\subsection{INT 15h, AX=2401h -- A20 Hattını Açma}

A20 adres hattı, 1\,MB üstündeki belleğe erişim için açılmalıdır (tarihsel bir
8086 uyumluluk kısıtı). BIOS üzerinden açmak en temiz yoldur.

\begin{ornekkutu}[A20 hattını BIOS ile açma ve doğrulama]
\begin{lstlisting}
    mov ax, 0x2403    ; once A20 destegi var mi diye sor
    int 0x15
    jc  .desteklenmiyor

    mov ax, 0x2402    ; A20 mevcut durumunu oku
    int 0x15
    cmp al, 1
    je  .zaten_acik   ; AL=1 ise zaten aciktir

    mov ax, 0x2401    ; A20 hattini ETKINLESTIR
    int 0x15
    jc  .hata
.zaten_acik:
.desteklenmiyor:
.hata:
\end{lstlisting}
\end{ornekkutu}

\begin{bilgikutu}
A20 açık değilse, 1\,MB sınırının üstündeki adresler "sarılır" (wrap-around) ve
0'dan başlar gibi davranır. Korumalı moda geçmeden \textbf{önce} A20'nin açık
olması şarttır; aksi halde çekirdeğiniz yüksek bellekte doğru çalışmaz. BIOS
yöntemi başarısız olursa alternatifler klavye denetleyicisi (port \hx{0x64})
veya "Fast A20" (port \hx{0x92}) üzerinden açmaktır.
\end{bilgikutu}

\subsection{INT 15h, AH=86h -- Mikrosaniye Cinsinden Bekleme}

Belirtilen süre kadar (\reg{CX:DX} mikrosaniye) programı bekletir. Gecikme
oluşturmanın hassas yoludur.

\begin{ornekkutu}[Tam olarak 500 milisaniye bekleme]
\begin{lstlisting}
    mov ah, 0x86
    mov cx, 0x0007    ; CX:DX = 500000 mikrosaniye
    mov dx, 0xA120    ;   (0x0007A120 = 500000)
    int 0x15          ; ~0.5 saniye bekle
\end{lstlisting}
\end{ornekkutu}

\subsection{INT 1Ah, AH=00h -- Sistem Saatini Okuma}

BIOS, saniyede yaklaşık 18.2 kez artan bir sayaç tutar. \reg{CX:DX} çiftinde
geri döner; bu değer basit zamanlama veya rastgele sayı tohumu (random seed)
için kullanılabilir.

\begin{ornekkutu}[Sistem tik sayacını okuyup gecikme oluşturma]
\begin{lstlisting}
; Baslangic zamanini oku
mov ah, 0x00
int 0x1A          ; CX:DX = gun basindan beri gecen tik sayisi
mov [baslangic], dx

bekle:
    mov ah, 0x00
    int 0x1A
    mov ax, dx
    sub ax, [baslangic]
    cmp ax, 18     ; ~18 tik = ~1 saniye
    jl  bekle      ; 1 saniye gecene kadar bekle
\end{lstlisting}
\end{ornekkutu}

\subsection{INT 12h -- Taban Bellek Boyutu}

Parametre almaz; \reg{AX}'te kullanılabilir taban bellek boyutunu KB cinsinden
döndürür (tipik olarak 640).

\begin{lstlisting}
int 0x12          ; AX = KB cinsinden taban bellek
; AX genellikle 640 (0x280) doner
\end{lstlisting}

% =====================================================================
\section{BIOS Örneklerinde Kullanılan Register'ların Açıklaması}
% =====================================================================

Bu bölüm, yukarıdaki BIOS kesme örneklerinde geçen her register'ın o çağrıda
\emph{neden} kullanıldığını tek tek açıklar. Böylece rehberin ilk yarısındaki
register bilgisi ile ikinci yarısındaki kesme örnekleri birbirine bağlanır.

\begin{bilgikutu}
\textbf{Hatırlatma -- genel desen:} Neredeyse tüm BIOS çağrılarında \reg{AH}
\emph{fonksiyon seçicisidir} (hangi alt işlemin çalışacağını belirler).
\reg{AL}, \reg{BX}, \reg{CX}, \reg{DX} \emph{parametreleri} taşır. Sonuç
register'larda ve/veya \reg{CF} (Carry) bayrağında döner. \reg{ES:BX},
\reg{DS:SI}, \reg{ES:DI} gibi \emph{çiftler} ise bir bellek adresini
(segment:offset) gösterir.
\end{bilgikutu}

% ---------------------------------------------------------------
\subsection{INT 10h, AH=0Eh -- Teletype Karakter Yazma}

\begin{center}
\begin{tabularx}{\textwidth}{@{}l l >{\raggedright\arraybackslash}X@{}}
\toprule
\rowcolor{anaMavi}\color{white}\textbf{Register} & \color{white}\textbf{Yön} & \color{white}\textbf{Bu örnekteki rolü} \\
\midrule
\reg{AH} & giriş & \hx{0x0E} yazılır -- "teletype yazma" fonksiyonunu seçer \\
\rowcolor{acikGri} \reg{AL} & giriş & Ekrana basılacak karakterin ASCII kodu (\kod{lodsb} her seferinde buraya bir bayt yükler) \\
\reg{SI} & giriş & Metnin başlangıç adresi; \kod{lodsb} bunu okuyup otomatik artırır \\
\rowcolor{acikGri} \reg{BH} & giriş & Görüntü sayfası (0); genelde önemsizdir \\
\bottomrule
\end{tabularx}
\end{center}

\begin{notkutu}
Örnekte \kod{lodsb} komutu şu üç işi tek adımda yapar: \reg{DS:SI}'nin
gösterdiği baytı \reg{AL}'ye yükler, sonra \reg{SI}'yi bir artırır. Yani
\reg{AL} ve \reg{SI} birlikte çalışır: \reg{SI} "nerede olduğumuzu", \reg{AL}
"o an hangi karakteri yazdığımızı" tutar. \kod{test al, al} ile \reg{AL}=0
(metin sonu) kontrol edilir.
\end{notkutu}

% ---------------------------------------------------------------
\subsection{INT 10h, AH=00h -- Video Modu Ayarlama}

\begin{center}
\begin{tabularx}{\textwidth}{@{}l l >{\raggedright\arraybackslash}X@{}}
\toprule
\rowcolor{anaMavi}\color{white}\textbf{Register} & \color{white}\textbf{Yön} & \color{white}\textbf{Bu örnekteki rolü} \\
\midrule
\reg{AH} & giriş & \hx{0x00} -- "video modu ayarla" fonksiyonu \\
\rowcolor{acikGri} \reg{AL} & giriş & İstenen mod numarası (örn. \hx{0x13} = 320$\times$200, 256 renk) \\
\bottomrule
\end{tabularx}
\end{center}
Burada yalnızca \reg{AX} kullanılır: \reg{AH} ne yapılacağını, \reg{AL} hangi
modun seçileceğini söyler. Çağrı ekranı da temizler.

% ---------------------------------------------------------------
\subsection{INT 10h, AH=02h -- İmleç Konumu}

\begin{center}
\begin{tabularx}{\textwidth}{@{}l l >{\raggedright\arraybackslash}X@{}}
\toprule
\rowcolor{anaMavi}\color{white}\textbf{Register} & \color{white}\textbf{Yön} & \color{white}\textbf{Bu örnekteki rolü} \\
\midrule
\reg{AH} & giriş & \hx{0x02} -- "imleç konumu ayarla" \\
\rowcolor{acikGri} \reg{BH} & giriş & Görüntü sayfası (0) \\
\reg{DH} & giriş & Hedef satır (0--24) -- \reg{DX}'in üst baytı \\
\rowcolor{acikGri} \reg{DL} & giriş & Hedef sütun (0--79) -- \reg{DX}'in alt baytı \\
\bottomrule
\end{tabularx}
\end{center}

\begin{ipucukutu}
Burada \reg{DH} ve \reg{DL}'nin ayrı ayrı kullanılması, tek bir 16-bit
register'ın (\reg{DX}) iki 8-bit parçaya bölünmesinin klasik örneğidir:
satır ve sütun tek bir register'da paketlenmiştir. \reg{DH}=satır, \reg{DL}=sütun.
\end{ipucukutu}

% ---------------------------------------------------------------
\subsection{INT 10h, AH=09h -- Renkli/Tekrarlı Karakter}

\begin{center}
\begin{tabularx}{\textwidth}{@{}l l >{\raggedright\arraybackslash}X@{}}
\toprule
\rowcolor{anaMavi}\color{white}\textbf{Register} & \color{white}\textbf{Yön} & \color{white}\textbf{Bu örnekteki rolü} \\
\midrule
\reg{AH} & giriş & \hx{0x09} -- "renkli karakter yaz" \\
\rowcolor{acikGri} \reg{AL} & giriş & Yazılacak karakter (\kod{'*'}) \\
\reg{BH} & giriş & Görüntü sayfası (0) \\
\rowcolor{acikGri} \reg{BL} & giriş & Renk özniteliği: üst 4 bit arka plan, alt 4 bit ön plan (\hx{0xF4} = beyaz zemin/kırmızı) \\
\reg{CX} & giriş & Karakterin kaç kez tekrarlanacağı (10) \\
\bottomrule
\end{tabularx}
\end{center}
Burada \reg{CX}'in "sayaç" rolü belirgindir: kaç kopya yazılacağını o tutar.
\reg{BL} ise tek bir baytta iki renk bilgisini (ön/arka) paketler.

% ---------------------------------------------------------------
\subsection{INT 10h, AH=06h -- Ekranı Kaydır/Temizle}

\begin{center}
\begin{tabularx}{\textwidth}{@{}l l >{\raggedright\arraybackslash}X@{}}
\toprule
\rowcolor{anaMavi}\color{white}\textbf{Register} & \color{white}\textbf{Yön} & \color{white}\textbf{Bu örnekteki rolü} \\
\midrule
\reg{AH} & giriş & \hx{0x06} -- "yukarı kaydır / temizle" \\
\rowcolor{acikGri} \reg{AL} & giriş & Kaydırılacak satır sayısı; \hx{0} = tüm pencereyi temizle \\
\reg{BH} & giriş & Yeni boş alanın renk özniteliği (\hx{0x1F} = mavi/beyaz) \\
\rowcolor{acikGri} \reg{CH}/\reg{CL} & giriş & \reg{CX} = pencerenin sol-üst köşesi (satır/sütun) \\
\reg{DH}/\reg{DL} & giriş & \reg{DX} = pencerenin sağ-alt köşesi (\hx{0x184F} = 24/79) \\
\bottomrule
\end{tabularx}
\end{center}
Burada \reg{CX} ve \reg{DX}, her biri iki 8-bitlik parçaya bölünerek bir
dikdörtgenin iki köşesini (4 koordinat) toplam iki register'da taşır.

% ---------------------------------------------------------------
\subsection{INT 10h, AH=0Ch -- Piksel Çizme}

\begin{center}
\begin{tabularx}{\textwidth}{@{}l l >{\raggedright\arraybackslash}X@{}}
\toprule
\rowcolor{anaMavi}\color{white}\textbf{Register} & \color{white}\textbf{Yön} & \color{white}\textbf{Bu örnekteki rolü} \\
\midrule
\reg{AH} & giriş & \hx{0x0C} -- "piksel yaz" \\
\rowcolor{acikGri} \reg{AL} & giriş & Pikselin renk numarası (0--255) \\
\reg{CX} & giriş & X koordinatı (sütun) -- döngüde \kod{inc cx} ile artırılır \\
\rowcolor{acikGri} \reg{DX} & giriş & Y koordinatı (satır) -- sabit tutulur \\
\bottomrule
\end{tabularx}
\end{center}
Çizgi çizme örneğinde \reg{DX} (Y) sabit kalırken \reg{CX} (X) her turda
artırılır; böylece yatay bir çizgi oluşur. Bu, register'ların bir döngü
değişkeni olarak nasıl kullanıldığını gösterir.

% ---------------------------------------------------------------
\subsection{INT 16h, AH=00h ve AH=01h -- Klavye}

\begin{center}
\begin{tabularx}{\textwidth}{@{}l l >{\raggedright\arraybackslash}X@{}}
\toprule
\rowcolor{anaMavi}\color{white}\textbf{Register} & \color{white}\textbf{Yön} & \color{white}\textbf{Rolü} \\
\midrule
\reg{AH} & giriş & \hx{0x00} = tuş bekle, \hx{0x01} = tuş var mı kontrol et \\
\rowcolor{acikGri} \reg{AL} & \textbf{çıkış} & Basılan tuşun ASCII kodu (BIOS doldurur) \\
\reg{AH} & \textbf{çıkış} & Tuşun tarama kodu (scan code) \\
\rowcolor{acikGri} \reg{ZF} & \textbf{çıkış} & (Yalnızca AH=01h) tampon boşsa 1, doluysa 0 \\
\bottomrule
\end{tabularx}
\end{center}

\begin{notkutu}
Bu örnek, register'ların hem \emph{giriş} hem \emph{çıkış} taşıyabildiğini
gösterir: çağrıdan önce \reg{AH}'ye fonksiyon numarasını yazarız, çağrıdan
sonra \emph{aynı} \reg{AX} register'ının içinde BIOS bize sonucu (\reg{AL}=ASCII,
\reg{AH}=scancode) geri verir. AH=01h'de asıl sonuç \reg{ZF} bayrağındadır --
bu yüzden çağrıdan hemen sonra \kod{jz} ile kontrol edilir.
\end{notkutu}

% ---------------------------------------------------------------
\subsection{INT 13h, AH=02h -- Diskten Sektör Okuma}

Bu, en çok register kullanan örnektir. Her biri CHS adreslemesinin bir
parçasını taşır:

\begin{center}
\small
\begin{tabularx}{\textwidth}{@{}l l >{\raggedright\arraybackslash}X@{}}
\toprule
\rowcolor{anaMavi}\color{white}\textbf{Register} & \color{white}\textbf{Yön} & \color{white}\textbf{Bu örnekteki rolü} \\
\midrule
\reg{AH} & giriş & \hx{0x02} -- "sektör oku" fonksiyonu \\
\rowcolor{acikGri} \reg{AL} & giriş & Kaç sektör okunacağı (1) \\
\reg{CH} & giriş & Silindir (cylinder) numarası \\
\rowcolor{acikGri} \reg{CL} & giriş & Sektör numarası (\textbf{1'den başlar!}) \\
\reg{DH} & giriş & Kafa (head) numarası \\
\rowcolor{acikGri} \reg{DL} & giriş & Sürücü numarası (\hx{0x80} = ilk sabit disk) \\
\reg{ES:BX} & giriş & Okunan verinin yükleneceği bellek adresi (segment:offset) \\
\rowcolor{acikGri} \reg{CF} & \textbf{çıkış} & 0 = başarılı, 1 = hata \\
\reg{AH} & \textbf{çıkış} & Hata durumunda hata kodu \\
\reg{DI} & yerel & Örnekteki yeniden deneme sayacı (BIOS kullanmaz, biz kullanırız) \\
\bottomrule
\end{tabularx}
\end{center}

\begin{ipucukutu}
Dikkat: \reg{CX} register'ı burada tek başına bir sayı değil, \emph{iki ayrı
bilgiye} bölünmüştür -- \reg{CH} silindiri, \reg{CL} sektörü taşır. Benzer
şekilde \reg{DX} de \reg{DH} (kafa) ve \reg{DL} (sürücü) olarak ikiye ayrılır.
\reg{ES:BX} çifti ise "veriyi belleğin neresine koyayım?" sorusunu yanıtlar.
Örnekteki \reg{DI} ise BIOS'a ait değildir; biz onu yeniden deneme sayacı
olarak kullanırız.
\end{ipucukutu}

% ---------------------------------------------------------------
\subsection{INT 13h, AH=42h -- LBA Okuma}

\begin{center}
\begin{tabularx}{\textwidth}{@{}l l >{\raggedright\arraybackslash}X@{}}
\toprule
\rowcolor{anaMavi}\color{white}\textbf{Register} & \color{white}\textbf{Yön} & \color{white}\textbf{Rolü} \\
\midrule
\reg{AH} & giriş & \hx{0x42} -- "genişletilmiş okuma" \\
\rowcolor{acikGri} \reg{DL} & giriş & Sürücü numarası (\hx{0x80}) \\
\reg{DS:SI} & giriş & Disk Adres Paketi'nin (DAP) adresi \\
\rowcolor{acikGri} \reg{CF} & \textbf{çıkış} & 0 = başarılı, 1 = hata \\
\bottomrule
\end{tabularx}
\end{center}
CHS'nin aksine burada tüm parametreler (sektör sayısı, hedef, LBA numarası)
\reg{DS:SI}'nin gösterdiği bellekteki bir \emph{yapıda} durur; register'lar
sadece o yapıyı işaret eder. Bu, çok sayıda parametreyi az register'la
geçirmenin yaygın yöntemidir.

% ---------------------------------------------------------------
\subsection{INT 15h, EAX=E820h -- Bellek Haritası}

\begin{center}
\small
\begin{tabularx}{\textwidth}{@{}l l >{\raggedright\arraybackslash}X@{}}
\toprule
\rowcolor{anaMavi}\color{white}\textbf{Register} & \color{white}\textbf{Yön} & \color{white}\textbf{Bu örnekteki rolü} \\
\midrule
\reg{EAX} & giriş/çıkış & Girişte \hx{0xE820}; her çağrıda yeniden ayarlanır. Çıkışta \kod{'SMAP'} imzası döner \\
\rowcolor{acikGri} \reg{EBX} & giriş/çıkış & "Devam" göstergesi: girişte 0, her çağrıda bir sonraki bölge için güncellenir; 0 olunca liste biter \\
\reg{ECX} & giriş & Bir girdinin bayt cinsinden boyutu (24) \\
\rowcolor{acikGri} \reg{EDX} & giriş & Sabit \hx{0x534D4150} imzası (\kod{'SMAP'}) \\
\reg{ES:DI} & giriş & Bölge tanımının yazılacağı tampon adresi \\
\rowcolor{acikGri} \reg{CF} & \textbf{çıkış} & 1 = hata veya listenin sonu \\
\bottomrule
\end{tabularx}
\end{center}

\begin{notkutu}
Bu, \reg{EBX}'in "durum taşıyıcı" olarak kullanıldığı güzel bir örnektir:
BIOS \reg{EBX}'e her seferinde kaldığı yeri yazar, biz de onu değiştirmeden
bir sonraki çağrıya veririz. Yani \reg{EBX} çağrılar arasında \emph{iterator}
(yineleyici) görevi görür. \reg{EAX} ve \reg{ECX}'in her turda yeniden
ayarlanması gerekir çünkü BIOS bunların üzerine yazar.
\end{notkutu}

% ---------------------------------------------------------------
\subsection{Diğer INT 15h Fonksiyonları}

\begin{center}
\begin{tabularx}{\textwidth}{@{}l l >{\raggedright\arraybackslash}X@{}}
\toprule
\rowcolor{anaMavi}\color{white}\textbf{Fonksiyon} & \color{white}\textbf{Register} & \color{white}\textbf{Rolü} \\
\midrule
E801h & \reg{AX}=\hx{0xE801} & Fonksiyon seçici \\
\rowcolor{acikGri} & \reg{AX}/\reg{CX} (çıkış) & 1--16\,MB arası bellek (KB) \\
& \reg{BX}/\reg{DX} (çıkış) & 16\,MB üstü bellek (64\,KB birim) \\
\midrule
88h & \reg{AH}=\hx{0x88} & Fonksiyon seçici \\
\rowcolor{acikGri} & \reg{AX} (çıkış) & 1\,MB üstü bellek (KB) \\
\midrule
2401h & \reg{AX}=\hx{0x2401} & A20'yi aç; \reg{CF} çıkışta hata durumu \\
\midrule
\rowcolor{acikGri} 86h & \reg{AH}=\hx{0x86} & Bekleme fonksiyonu \\
& \reg{CX:DX} & Mikrosaniye cinsinden süre (32-bit değer iki register'a bölünür) \\
\bottomrule
\end{tabularx}
\end{center}

\begin{ipucukutu}
\reg{CX:DX} çiftinde 32-bit bir sayının nasıl taşındığına dikkat edin:
\reg{CX} üst 16 biti, \reg{DX} alt 16 biti tutar. 16-bit modda 32-bit değer
tek register'a sığmadığı için bu "register çifti" tekniği kullanılır --
tıpkı çarpma sonucunun \reg{DX:AX}'te dönmesi gibi.
\end{ipucukutu}

% ---------------------------------------------------------------
\subsection{INT 1Ah (Saat) ve INT 12h (Bellek)}

\begin{center}
\begin{tabularx}{\textwidth}{@{}l l >{\raggedright\arraybackslash}X@{}}
\toprule
\rowcolor{anaMavi}\color{white}\textbf{Kesme} & \color{white}\textbf{Register} & \color{white}\textbf{Rolü} \\
\midrule
INT 1Ah & \reg{AH}=\hx{0x00} & Fonksiyon: sayacı oku \\
\rowcolor{acikGri} & \reg{CX:DX} (çıkış) & Gün başından beri geçen tik sayısı (32-bit, çift register) \\
\midrule
INT 12h & (giriş yok) & Parametre almaz \\
\rowcolor{acikGri} & \reg{AX} (çıkış) & Taban bellek boyutu (KB) \\
\bottomrule
\end{tabularx}
\end{center}
INT 12h, hiç giriş parametresi almayan nadir bir kesmedir; sonucu doğrudan
\reg{AX}'te verir. INT 1Ah ise saat değerini yine bir \reg{CX:DX} çiftinde
döndürür.

% ---------------------------------------------------------------
\subsection{Özet: Register'ların BIOS Çağrılarındaki Rolleri}

\begin{center}
\begin{tabularx}{\textwidth}{@{}l >{\raggedright\arraybackslash}X@{}}
\toprule
\rowcolor{anaMavi}\color{white}\textbf{Register} & \color{white}\textbf{BIOS çağrılarındaki tipik rolü} \\
\midrule
\reg{AH} & Neredeyse her zaman \textbf{fonksiyon seçici} (hangi alt işlem) \\
\rowcolor{acikGri} \reg{AL} & Karakter, mod numarası, sektör sayısı gibi tekil parametre; sık sık çıkışta da kullanılır \\
\reg{BX} / \reg{BH} / \reg{BL} & Sayfa, renk özniteliği veya hedef offset (\reg{ES:BX}) \\
\rowcolor{acikGri} \reg{CX} / \reg{CH} / \reg{CL} & Sayaç, tekrar sayısı veya CHS'de silindir+sektör \\
\reg{DX} / \reg{DH} / \reg{DL} & Sürücü, koordinat, kafa veya \reg{CX:DX}'te 32-bit değerin alt yarısı \\
\rowcolor{acikGri} \reg{SI} / \reg{DI} & Bellek işaretçisi (metin, DAP, tampon adresi) \\
\reg{ES} / \reg{DS} & \reg{BX}/\reg{SI}/\reg{DI} ile birlikte adresin segment kısmı \\
\rowcolor{acikGri} \reg{CF} (bayrak) & Çıkışta \textbf{hata göstergesi} (1 = hata) \\
\reg{ZF} (bayrak) & Bazı çağrılarda durum çıkışı (örn. klavye tamponu boş mu) \\
\bottomrule
\end{tabularx}
\end{center}

% =====================================================================
\section{Tam Örnek: İki Aşamalı Önyükleyici}
% =====================================================================

Aşağıdaki tam çalışan örnek, bir önyükleyicinin tipik görevlerini bir arada
gösterir: segmentleri kur, mesaj yazdır, diskten ikinci aşamayı yükle ve ona
atla. Bu, gerçek işletim sistemlerinin önyükleme mantığının basitleştirilmiş
hâlidir.

\begin{lstlisting}
bits 16               ; 16-bit gercek mod
org 0x7C00            ; BIOS onyukleyiciyi buraya yukler

start:
    ; --- 1. Segmentleri ve yigini guvenli sekilde ayarla ---
    cli               ; kesmeleri kapat (kurulum sirasinda)
    xor ax, ax
    mov ds, ax        ; DS = 0
    mov es, ax        ; ES = 0
    mov ss, ax        ; SS = 0
    mov sp, 0x7C00    ; yigin onyukleyicinin altinda
    sti               ; kesmeleri geri ac

    ; --- 2. Karsilama mesaji yazdir ---
    mov si, mesaj
    call yaz_metin

    ; --- 3. Diskten ikinci asamayi yukle (sektor 2) ---
    mov di, 3         ; 3 deneme hakki
.oku:
    mov ah, 0x02      ; sektor oku
    mov al, 1         ; 1 sektor
    mov ch, 0         ; silindir 0
    mov cl, 2         ; sektor 2
    mov dh, 0         ; kafa 0
    mov dl, 0x80      ; ilk sabit disk
    mov bx, 0x8000    ; ES:BX = 0000:8000 hedef adres
    int 0x13
    jnc .yuklendi
    xor ah, ah        ; hata -> diski resetle
    int 0x13
    dec di
    jnz .oku
    mov si, hata_msj  ; 3 deneme basarisiz -> hata mesaji
    call yaz_metin
    jmp dur

.yuklendi:
    ; --- 4. Ikinci asamaya atla ---
    jmp 0x0000:0x8000 ; 0x8000 adresindeki koda gec

dur:
    hlt
    jmp dur

; --- Yardimci rutin: SI'daki null-terminated metni yazar ---
yaz_metin:
    mov ah, 0x0E
.dongu:
    lodsb
    test al, al
    jz  .son
    int 0x10
    jmp .dongu
.son:
    ret

mesaj:    db "Onyukleyici basladi, kernel yukleniyor...", 0x0D, 0x0A, 0
hata_msj: db "DISK HATASI!", 0x0D, 0x0A, 0

times 510-($-$$) db 0 ; 510. bayta kadar sifirla doldur
dw 0xAA55             ; onyukleme imzasi (son 2 bayt)
\end{lstlisting}

\begin{notkutu}[Önyükleme sektörü kuralları]
Önyükleme sektörü \textbf{tam 512 bayt} olmalı ve son 2 baytı \hx{0xAA55}
imzasını içermelidir; BIOS bu imzayı görmezse diski önyüklenebilir kabul etmez.
\kod{times 510-(\$-\$\$) db 0} ifadesi, kodun bittiği yerden 510. bayta kadar
olan boşluğu sıfırla doldurur (\kod{\$} mevcut konum, \kod{\$\$} bölüm başıdır).
\end{notkutu}

% =====================================================================
\section{Korumalı Moda Geçiş}
% =====================================================================

Önyükleyicinin son görevi, 16-bit gerçek moddan 32-bit korumalı moda geçmektir.
Bunun için bir GDT (Global Descriptor Table) hazırlanır ve \reg{CR0}'ın PE
biti set edilir.

\begin{ornekkutu}[Gerçek moddan korumalı moda geçiş]
\begin{lstlisting}
    cli               ; kesmeleri kapat (IDT henuz hazir degil)
    lgdt [gdt_tanim]  ; GDT'yi islemciye tanit

    mov eax, cr0
    or  eax, 1        ; PE bitini (bit 0) set et
    mov cr0, eax      ; artik korumali moddayiz!

    ; boru hattini temizlemek icin uzak atlama sart:
    jmp 0x08:korumali_mod  ; 0x08 = GDT'deki kod segmenti secici

bits 32               ; artik 32-bit kod
korumali_mod:
    mov ax, 0x10      ; 0x10 = veri segmenti secici
    mov ds, ax
    mov es, ax
    mov ss, ax
    mov esp, 0x90000  ; yeni 32-bit yigin
    ; ... 32-bit cekirdek kodu buradan devam eder ...

; --- Global Descriptor Table ---
gdt_baslangic:
    dq 0x0000000000000000  ; null tanimlayici (zorunlu)
gdt_kod:
    dq 0x00CF9A000000FFFF  ; kod segmenti (taban=0, limit=4GB)
gdt_veri:
    dq 0x00CF92000000FFFF  ; veri segmenti (taban=0, limit=4GB)
gdt_son:

gdt_tanim:
    dw gdt_son - gdt_baslangic - 1  ; GDT boyutu - 1
    dd gdt_baslangic                ; GDT adresi
\end{lstlisting}
\end{ornekkutu}

\begin{uyarikutu}
\reg{CR0}'ın PE bitini set ettikten hemen sonra \textbf{uzak atlama} (far jump)
yapmak zorunludur. Bu, işlemcinin komut boru hattını (pipeline) temizler ve
\reg{CS} register'ını yeni 32-bit segment seçicisiyle yükler. Bu adım
atlanırsa işlemci tanımsız duruma girer.
\end{uyarikutu}

% =====================================================================
\section{Sayfalama (Paging)}
% =====================================================================

Sayfalama, sanal adresleri fiziksel adreslere çeviren donanım mekanizmasıdır.
Bellek 4\,KB'lık "sayfa"lara bölünür; bir sayfa tablosu her sanal sayfayı bir
fiziksel çerçeveye eşler. Bu sayede her programa kendi izole adres alanı verilir.

\subsection{Sayfa Tablosu Hiyerarşisi}

\begin{center}
\begin{tabularx}{\textwidth}{@{}l l >{\raggedright\arraybackslash}X@{}}
\toprule
\rowcolor{anaMavi}\color{white}\textbf{Mod} & \color{white}\textbf{Seviye} & \color{white}\textbf{Yapı} \\
\midrule
32-bit & 2 seviye & Sayfa Dizini (PD) $\rightarrow$ Sayfa Tablosu (PT) \\
\rowcolor{acikGri} 32-bit PAE & 3 seviye & PDPT $\rightarrow$ PD $\rightarrow$ PT \\
64-bit & 4 seviye & PML4 $\rightarrow$ PDPT $\rightarrow$ PD $\rightarrow$ PT \\
\bottomrule
\end{tabularx}
\end{center}

\begin{bilgikutu}
Bir sanal adres, tablo seviyelerine bölünerek çözümlenir. 32-bit'te 32 bitlik
adres şöyle ayrışır: üst 10 bit sayfa dizini indeksi, ortadaki 10 bit sayfa
tablosu indeksi, alt 12 bit ise sayfa içi offset (4\,KB = $2^{12}$). İşlemci
bu çeviriyi otomatik yapar; TLB (Translation Lookaside Buffer) sık kullanılan
çevirileri önbelleğe alarak hızlandırır.
\end{bilgikutu}

\subsection{32-bit Sayfalamayı Etkinleştirme}

Aşağıdaki örnek, ilk 4\,MB'lık belleği birebir (identity mapping) eşler --
yani sanal adres = fiziksel adres. Çekirdeğin kendi kodunu çalıştırmaya devam
edebilmesi için bu şarttır.

\begin{ornekkutu}[Birebir eşlemeli sayfa tablosu kurma (32-bit)]
\begin{lstlisting}
bits 32
kur_sayfalama:
    ; --- 1. Sayfa tablosunu birebir eslemeyle doldur ---
    mov edi, 0x101000     ; sayfa tablosunun adresi (PT)
    mov cr3, edi          ; ileride PD adresini yazacagiz
    mov eax, 0x00000003   ; ilk sayfa: adres=0, bayrak=Present+RW
    mov ecx, 1024         ; 1024 girdi = 4MB kapsam
.doldur:
    mov [edi], eax        ; girdiyi yaz
    add eax, 0x1000       ; sonraki 4KB sayfaya gec
    add edi, 4            ; sonraki girdi (her girdi 4 bayt)
    loop .doldur

    ; --- 2. Sayfa dizinini kur (tek girdi PT'yi gosterir) ---
    mov edi, 0x100000     ; sayfa dizini (PD) adresi
    mov eax, 0x101003     ; PT adresi | Present | RW
    mov [edi], eax
    mov dword [edi+4], 0  ; diger PD girdileri bos (Present=0)

    ; --- 3. CR3'e sayfa dizini adresini yukle ---
    mov eax, 0x100000
    mov cr3, eax

    ; --- 4. CR0'in PG bitini (bit 31) set ederek sayfalamayi ac ---
    mov eax, cr0
    or  eax, 0x80000000
    mov cr0, eax          ; sayfalama artik AKTIF
    ret
\end{lstlisting}
\end{ornekkutu}

\subsection{Sayfa Tablosu Girdisi (PTE) Bayrakları}

Her sayfa tablosu girdisinin alt 12 biti kontrol bayraklarıdır (adresin alt 12
biti zaten 0 olduğu için bu bitler boştur ve bayrak olarak kullanılır).

\begin{center}
\begin{tabularx}{\textwidth}{@{}c l >{\raggedright\arraybackslash}X@{}}
\toprule
\rowcolor{anaMavi}\color{white}\textbf{Bit} & \color{white}\textbf{Ad} & \color{white}\textbf{Anlamı} \\
\midrule
0 & P (Present) & Sayfa bellekte mevcut mu? 0 ise erişim page fault üretir \\
\rowcolor{acikGri} 1 & R/W & 1=yazılabilir, 0=salt okunur \\
2 & U/S & 1=kullanıcı (ring 3), 0=yalnızca çekirdek (ring 0) \\
\rowcolor{acikGri} 3 & PWT & Write-through önbellekleme \\
4 & PCD & Önbelleği devre dışı bırak \\
\rowcolor{acikGri} 5 & A (Accessed) & İşlemci erişince otomatik set edilir \\
6 & D (Dirty) & Sayfaya yazılınca set edilir (swap için kullanılır) \\
\rowcolor{acikGri} 63 & NX & (64-bit) Bu sayfada kod çalıştırmayı yasakla \\
\bottomrule
\end{tabularx}
\end{center}

\begin{uyarikutu}
Sayfalamayı açtıktan sonra, o an çalışan kodun bulunduğu adres birebir eşlenmiş
\textbf{olmalıdır}. Aksi halde \reg{CR0}'a yazan komuttan sonraki komutu getirmeye
çalışan işlemci geçersiz bir sanal adrese bakar ve tripla hata (triple fault)
ile yeniden başlar. Bu yüzden önce identity mapping yapılır.
\end{uyarikutu}

\begin{ipucukutu}
Bir sayfa tablosu girdisini değiştirdiğinizde TLB'deki eski çeviri hâlâ
geçerli olabilir. \kod{invlpg [adres]} komutuyla ilgili TLB girdisini
geçersizleştirebilir, ya da \reg{CR3}'ü yeniden yükleyerek tüm TLB'yi
temizleyebilirsiniz.
\end{ipucukutu}

% =====================================================================
\section{Kesme İşleyici (ISR) Yazma}
% =====================================================================

Şimdiye kadar BIOS'un \emph{sağladığı} kesmeleri kullandık. Bir işletim sistemi
ise kendi kesme işleyicilerini (ISR -- Interrupt Service Routine) yazar.
Korumalı modda bunun için \textbf{IDT} (Interrupt Descriptor Table) kurulur.

\subsection{IDT -- Kesme Tanımlayıcı Tablosu}

Gerçek moddaki basit IVT'nin yerini korumalı modda IDT alır. IDT'de her girdi
8 bayttır ve bir kesme işleyicisinin adresini, segment seçicisini ve tipini
tutar.

\begin{center}
\begin{tabularx}{\textwidth}{@{}l >{\raggedright\arraybackslash}X@{}}
\toprule
\rowcolor{acikMavi}\color{anaMavi}\textbf{Alan} & \color{anaMavi}\textbf{Açıklama} \\
\midrule
Offset (0--15) & İşleyici adresinin alt 16 biti \\
Seçici & Kod segmenti seçicisi (genellikle \hx{0x08}) \\
Tip \& öznitelik & Kapı tipi: \hx{0x8E}=kesme kapısı (ring 0) \\
Offset (16--31) & İşleyici adresinin üst 16 biti \\
\bottomrule
\end{tabularx}
\end{center}

\subsection{Bir IDT Girdisi Ayarlama}

\begin{ornekkutu}[Belirli bir kesme numarasına işleyici atama rutini]
\begin{lstlisting}
bits 32
; Giris: AL = kesme numarasi, EBX = isleyici adresi
idt_girdi_ayarla:
    push edi
    movzx edi, al         ; EDI = kesme numarasi
    shl edi, 3            ; her girdi 8 bayt -> indeks*8
    add edi, idt          ; IDT icindeki girdi adresi

    mov [edi],   bx       ; offset bit 0-15
    mov word [edi+2], 0x08 ; kod segmenti secicisi
    mov byte [edi+4], 0    ; ayrilmis (0)
    mov byte [edi+5], 0x8E ; tip: 32-bit kesme kapisi, ring 0
    shr ebx, 16
    mov [edi+6], bx       ; offset bit 16-31
    pop edi
    ret

; IDT: 256 girdi x 8 bayt
idt:      times 256*8 db 0
idt_tanim:
    dw 256*8 - 1          ; IDT boyutu - 1
    dd idt                ; IDT adresi
\end{lstlisting}
\end{ornekkutu}

\subsection{Basit Bir Kesme İşleyicisi}

İşleyiciler \kod{iret} (interrupt return) ile bitmelidir; \kod{ret} değil.
İşleyici, kullandığı tüm register'ları koruyup geri yüklemelidir.

\begin{ornekkutu}[Ekrana nokta basan basit ISR (örneğin zamanlayıcı)]
\begin{lstlisting}
bits 32
zamanlayici_isr:
    pusha                 ; TUM genel register'lari sakla

    ; --- isleyici govdesi: ekranin sol ustune '.' yaz ---
    mov edi, 0xB8000      ; VGA metin belleginin adresi
    mov byte [edi], '.'   ; karakter
    mov byte [edi+1], 0x0A ; renk: acik yesil

    ; --- PIC'e "kesme islendi" bildir (End Of Interrupt) ---
    mov al, 0x20          ; EOI komutu
    out 0x20, al          ; ana PIC'e gonder (port 0x20)

    popa                  ; register'lari geri yukle
    iret                  ; kesmeden don (EFLAGS + CS:EIP geri yuklenir)
\end{lstlisting}
\end{ornekkutu}

\begin{uyarikutu}
Bir donanım kesmesi işleyicisinin sonunda \textbf{PIC'e EOI (End Of Interrupt)}
sinyali göndermek zorunludur (\kod{out 0x20, al} ile \hx{0x20} değeri). Aksi
halde PIC o kesme hattının daha fazla kesme göndermesini engeller ve sistem
kilitlenir. Slave PIC'ten (IRQ 8--15) gelen kesmelerde hem \hx{0xA0} hem de
\hx{0x20} portlarına EOI gönderilir.
\end{uyarikutu}

\subsection{PIC'i Yeniden Haritalama}

Varsayılan olarak donanım kesmeleri (IRQ 0--15), işlemcinin istisna
(exception) numaralarıyla (0--31) çakışır. Bu yüzden PIC (8259A denetleyici)
yeniden programlanarak IRQ'lar \hx{0x20}--\hx{0x2F} aralığına taşınır.

\begin{ornekkutu}[PIC'i 0x20 ofsetine yeniden haritalama]
\begin{lstlisting}
bits 32
pic_yeniden_haritala:
    mov al, 0x11          ; baslatma komutu (ICW1)
    out 0x20, al          ; ana PIC'e
    out 0xA0, al          ; slave PIC'e

    mov al, 0x20          ; ICW2: ana PIC ofseti -> IRQ0 = int 0x20
    out 0x21, al
    mov al, 0x28          ; slave PIC ofseti -> IRQ8 = int 0x28
    out 0xA1, al

    mov al, 0x04          ; ICW3: ana-slave baglantisi (IRQ2)
    out 0x21, al
    mov al, 0x02
    out 0xA1, al

    mov al, 0x01          ; ICW4: 8086 modu
    out 0x21, al
    out 0xA1, al

    mov al, 0x00          ; maskeyi temizle: tum kesmeler acik
    out 0x21, al
    out 0xA1, al
    ret
\end{lstlisting}
\end{ornekkutu}

\subsection{Klavye Kesme İşleyicisi (IRQ1)}

Klavye, IRQ1 (yeniden haritaladıktan sonra \hx{INT 0x21}) üretir. İşleyici,
tarama kodunu port \hx{0x60}'tan okumalıdır.

\begin{ornekkutu}[Klavyeden tarama kodu okuyan ISR]
\begin{lstlisting}
bits 32
klavye_isr:
    pusha
    in  al, 0x60          ; klavye veri portundan tarama kodunu oku

    ; AL simdi tarama kodunu icerir.
    ; Ust bit 1 ise tus BIRAKILDI, 0 ise BASILDI:
    test al, 0x80
    jnz .birakildi

    ; --- tusa basildi: burada isle ---
    ; ornek: tarama kodunu bir tampona kaydet
    mov [son_tus], al

.birakildi:
    mov al, 0x20          ; PIC'e EOI gonder
    out 0x20, al
    popa
    iret

son_tus: db 0
\end{lstlisting}
\end{ornekkutu}

\begin{bilgikutu}
Kurulum sırası önemlidir: (1) IDT'yi kur ve \kod{lidt} ile yükle, (2) PIC'i
yeniden haritala, (3) işleyicileri IDT'ye kaydet, (4) \kod{sti} ile kesmeleri
aç. Bu adımlardan biri atlanırsa ilk donanım kesmesinde sistem çöker. Klavye
tarama kodu ASCII \emph{değildir}; bir çeviri tablosuyla (scancode$\rightarrow$ASCII)
dönüştürülmelidir.
\end{bilgikutu}

% =====================================================================
\section{Hızlı Referans: Register Kullanım Kalıpları}
% =====================================================================

\begin{center}
\begin{tabularx}{\textwidth}{@{}>{\raggedright\arraybackslash}X l@{}}
\toprule
\rowcolor{anaMavi}\color{white}\textbf{Amaç} & \color{white}\textbf{İlgili Register'lar} \\
\midrule
Fonksiyon dönüş değeri & \reg{RAX}/\reg{EAX} \\
\rowcolor{acikGri} Döngü sayacı & \reg{RCX}/\reg{ECX} \\
Çarpma/bölme (yüksek bölüm) & \reg{RDX}/\reg{EDX} \\
\rowcolor{acikGri} String kaynak/hedef & \reg{RSI} / \reg{RDI} \\
Yığın yönetimi & \reg{RSP} / \reg{RBP} \\
\rowcolor{acikGri} BIOS kesme fonksiyon seçimi & \reg{AH} \\
Kesme dönüş hata durumu & \reg{CF} bayrağı \\
\rowcolor{acikGri} G/Ç port numarası & \reg{DX} \\
\bottomrule
\end{tabularx}
\end{center}

\subsection{System V AMD64 Çağrı Kuralı (Linux 64-bit)}

Modern 64-bit Linux'ta fonksiyon parametreleri şu sırayla register'lardan geçer:

\begin{center}
\begin{tabularx}{\textwidth}{@{}c c c c c c c@{}}
\toprule
\rowcolor{anaMavi}
\color{white}\textbf{1.} & \color{white}\textbf{2.} & \color{white}\textbf{3.} & \color{white}\textbf{4.} & \color{white}\textbf{5.} & \color{white}\textbf{6.} & \color{white}\textbf{Dönüş} \\
\midrule
\reg{RDI} & \reg{RSI} & \reg{RDX} & \reg{RCX} & \reg{R8} & \reg{R9} & \reg{RAX} \\
\bottomrule
\end{tabularx}
\end{center}

\begin{bilgikutu}
Linux \kod{syscall} kuralında ise 4. parametre \reg{RCX} yerine \reg{R10}
kullanır ve syscall numarası \reg{RAX}'e yazılır. Ayrıca \reg{RBX}, \reg{RBP},
\reg{R12}--\reg{R15} "çağrı boyunca korunan" (callee-saved) register'lardır;
bir fonksiyon bunları değiştirirse geri yüklemekle yükümlüdür.
\end{bilgikutu}

\begin{ornekkutu}[64-bit Linux'ta "Merhaba" yazan syscall örneği]
\begin{lstlisting}
section .data
    mesaj db "Merhaba", 10   ; 10 = yeni satir
    uzunluk equ $ - mesaj    ; metnin uzunlugu

section .text
    global _start
_start:
    mov rax, 1        ; sys_write cagri numarasi
    mov rdi, 1        ; dosya tanimlayici: stdout
    mov rsi, mesaj    ; tampon adresi
    mov rdx, uzunluk  ; yazilacak bayt sayisi
    syscall           ; cekirdegi cagir

    mov rax, 60       ; sys_exit cagri numarasi
    xor rdi, rdi      ; cikis kodu 0
    syscall
\end{lstlisting}
\end{ornekkutu}

% =====================================================================
\section{Pratik Kod Örnekleri}
% =====================================================================

Bu bölüm, yukarıdaki kavramları birleştiren sık kullanılan rutinleri içerir.

\subsection{String Uzunluğu (strlen)}

Null ile biten bir metnin uzunluğunu sayar -- C'deki \kod{strlen} karşılığı.

\begin{ornekkutu}[SI'daki metnin uzunluğunu CX'te döndürme]
\begin{lstlisting}
strlen:               ; giris: SI = metin,  cikis: CX = uzunluk
    xor cx, cx        ; sayac = 0
.dongu:
    mov al, [si]      ; siradaki karakter
    test al, al       ; null mu?
    jz  .son
    inc cx            ; uzunlugu artir
    inc si            ; sonraki karaktere gec
    jmp .dongu
.son:
    ret
\end{lstlisting}
\end{ornekkutu}

\subsection{İki String'i Karşılaştırma (strcmp)}

\begin{ornekkutu}[SCASB/CMPSB olmadan elle karşılaştırma]
\begin{lstlisting}
strcmp:               ; SI ve DI: iki metin.  Esitse ZF=1 doner
.dongu:
    mov al, [si]
    mov bl, [di]
    cmp al, bl
    jne .farkli       ; karakterler farkli -> esit degil
    test al, al       ; ikisi de null mu? (metin sonu)
    jz  .esit         ; evet -> metinler ayni
    inc si
    inc di
    jmp .dongu
.farkli:
    ; ZF=0 (cmp farkli oldugu icin)
    ret
.esit:
    ; ZF=1
    ret
\end{lstlisting}
\end{ornekkutu}

\subsection{Sayıyı Onlu (Decimal) Olarak Ekrana Yazma}

Bir tamsayıyı ekrana rakamlarıyla basmak, sayıyı 10'a bölerek kalanları ters
sırada toplamayı gerektirir. Klasik bir alçak seviye problemi.

\begin{ornekkutu}[AX'teki işaretsiz sayıyı ekrana yazdırma (gerçek mod)]
\begin{lstlisting}
yaz_sayi:             ; giris: AX = yazdirilacak sayi
    push bx
    push cx
    push dx
    mov cx, 0         ; kac rakam yiginda oldugunu say
    mov bx, 10        ; bolen (taban 10)
.bol:
    xor dx, dx        ; DX'i temizle (DX:AX / BX icin)
    div bx            ; AX = AX/10,  DX = kalan (bir rakam)
    add dl, '0'       ; rakami ASCII'ye cevir
    push dx           ; yigina at (ters sirada birikir)
    inc cx
    test ax, ax       ; sayi bitti mi?
    jnz .bol
.yaz:
    pop dx            ; rakamlari dogru sirada geri cek
    mov ah, 0x0E
    mov al, dl
    int 0x10          ; ekrana yaz
    loop .yaz
    pop dx
    pop cx
    pop bx
    ret
\end{lstlisting}
\end{ornekkutu}

\subsection{Baytı Onaltılık (Hex) Olarak Yazma}

\begin{ornekkutu}[AL'deki baytı iki hex rakamı olarak yazma]
\begin{lstlisting}
yaz_hex:              ; giris: AL = bayt
    push ax
    mov ah, al        ; kopyala
    shr al, 4         ; ust yarim-bayti al (yuksek nibble)
    call .nibble
    mov al, ah
    and al, 0x0F      ; alt yarim-bayt (dusuk nibble)
    call .nibble
    pop ax
    ret
.nibble:              ; AL'deki 0-15 degerini hex rakamina cevir
    cmp al, 10
    jl  .rakam
    add al, 'A' - 10  ; 10-15 -> 'A'-'F'
    jmp .yaz
.rakam:
    add al, '0'       ; 0-9 -> '0'-'9'
.yaz:
    mov ah, 0x0E
    int 0x10
    ret
\end{lstlisting}
\end{ornekkutu}

\subsection{Basit Gecikme Döngüsü (Busy-wait)}

\begin{ornekkutu}[İç içe döngüyle kaba gecikme]
\begin{lstlisting}
gecikme:
    mov cx, 0xFFFF    ; dis dongu sayaci
.dis:
    mov dx, 0xFFFF    ; ic dongu sayaci
.ic:
    dec dx
    jnz .ic           ; ic dongu bitene kadar
    dec cx
    jnz .dis          ; dis dongu bitene kadar
    ret
\end{lstlisting}
\end{ornekkutu}

\begin{uyarikutu}
Busy-wait gecikmeleri işlemci hızına bağlıdır ve her makinede farklı süre
alır; bu yüzden gerçek uygulamalarda \kod{INT 15h AH=86h} (mikrosaniye) veya
zamanlayıcı kesmesi tercih edilir. Yukarıdaki döngü yalnızca kaba testler
içindir.
\end{uyarikutu}

% =====================================================================
\section{Derleme ve Test}
% =====================================================================

\begin{ipucukutu}
\textbf{Önyükleyici örneklerini test etmek} için NASM ile derleyip QEMU'da
çalıştırabilirsiniz:
\begin{lstlisting}[numbers=none]
nasm -f bin onyukleyici.asm -o onyukleyici.bin
qemu-system-x86_64 -drive format=raw,file=onyukleyici.bin
\end{lstlisting}
\textbf{64-bit Linux programını} derleyip çalıştırmak için:
\begin{lstlisting}[numbers=none]
nasm -f elf64 program.asm -o program.o
ld program.o -o program
./program
\end{lstlisting}
Hata ayıklama için QEMU'ya \kod{-s -S} ekleyip GDB ile \kod{target remote
localhost:1234} bağlanabilirsiniz.
\end{ipucukutu}

% =====================================================================
\section{Kapanış Notları}
% =====================================================================

Bu rehber, x86 register mimarisinin temel taşlarını, adresleme modlarını, SIMD
birimlerini ve gerçek modda en çok kullanılan BIOS kesmelerini bol örnekle
kapsadı. Özetle:

\begin{itemize}[leftmargin=1.4em]
  \item \textbf{Register'lar} işlemcinin en hızlı çalışma belleğidir; her birinin
        geleneksel bir amacı ve boyut hiyerarşisi (RAX$\to$EAX$\to$AX$\to$AL) vardır.
  \item \textbf{Adresleme modları}, özellikle ölçekli indeks, dizilere ve
        yapılara verimli erişim sağlar.
  \item \textbf{Bayraklar} karar mekanizmasının kalbidir; \kod{CMP}+koşullu
        atlama ikilisi tüm dallanmaların temelidir.
  \item \textbf{BIOS kesmeleri}, gerçek modda donanıma açılan standart kapılardır;
        \reg{AH} fonksiyonu, diğer register'lar parametreleri, \reg{CF} ise
        hata durumunu taşır.
  \item \textbf{Önyükleyici} bu parçaları birleştirir: segment kurulumu, mesaj,
        disk okuma ve korumalı moda geçiş.
\end{itemize}

İkisini birlikte kavramak -- register'lar ve kesmeler -- önyükleyici yazmaktan
işletim sistemi geliştirmeye kadar tüm alçak seviye programlamanın temelidir.

\vspace{0.8cm}
\begin{center}
{\color{turuncu}\rule{0.5\textwidth}{1.5pt}}\\[6pt]
{\color{ortaGri}\textit{Rehber sonu -- İyi çalışmalar!}}
\end{center}

\end{document}
